
\documentclass[11pt,a4paper]{ctexart}
\usepackage{geometry}
\geometry{left=1.55cm,right=1.55cm,top=1.35cm,bottom=1.25cm,headsep=0.32cm,footskip=0.5cm}
\usepackage{fontspec}
\setmainfont{Times New Roman}
\setsansfont{Helvetica}
\usepackage{xcolor}
\usepackage{booktabs}
\usepackage{tabularx}
\usepackage{array}
\usepackage{longtable}
\usepackage{amssymb}
\usepackage{enumitem}
\usepackage{fancyhdr}
\usepackage[hidelinks]{hyperref}
\hypersetup{pdftitle={英语二小作文·一文一模版},pdfauthor={DSH 学习助手},pdfsubject={考研英语二 Part A 应用文}}
\xeCJKDeclareCharClass{CJK}{"2460 -> "247F}
\xeCJKDeclareCharClass{CJK}{"25A0 -> "25FF}
\definecolor{accent}{HTML}{1F3A5F}
\definecolor{hot}{HTML}{A23B2A}
\definecolor{gray}{HTML}{5B5B5B}
\definecolor{rule}{HTML}{C9D3DD}
\definecolor{soft}{HTML}{F3F6F9}
\definecolor{sand}{HTML}{FBF4E6}
\definecolor{sandline}{HTML}{DCCBA8}
\renewcommand{\familydefault}{\rmdefault}
\setlength{\parindent}{0pt}
\setlength{\parskip}{0pt}
\linespread{1.02}
\pagestyle{fancy}
\fancyhf{}
\renewcommand{\headrulewidth}{0pt}
\fancyhead[L]{\footnotesize\color{gray} 英语二小作文 · 一文一模版}
\fancyhead[R]{\footnotesize\color{gray} Part A · 应用文}
\fancyfoot[L]{\footnotesize\color{gray} 12 类 · 每类一整篇}
\fancyfoot[R]{\footnotesize\color{gray} \thepage}
\renewcommand{\arraystretch}{1.12}
\setlength{\LTpre}{0pt}\setlength{\LTpost}{0pt}

\begin{document}

\clearpage
\thispagestyle{empty}
\vspace*{0.55cm}
{\color{accent}\footnotesize\bfseries 考研英语（二）\quad Part A\quad 应用文写作}\par
\vspace{10pt}
{\color{accent}\rule{\linewidth}{2.4pt}}\par
\vspace{14pt}
{\fontsize{40}{44}\selectfont\bfseries 一文一模版}\par
\vspace{8pt}
{\LARGE\bfseries 英语二小作文 · 12 类文体 · 每类一整篇}\par
\vspace{8pt}
{\large\color{gray} 背一篇 → 换五个词 → 直接上考场}\par
\vspace{6pt}
{\color{accent}\rule{\linewidth}{0.9pt}}\par
\vspace{16pt}

\noindent\begin{minipage}[t]{0.32\linewidth}
{\color{accent}\fontsize{30}{32}\selectfont\bfseries 12}\par
{\small\color{gray} 类文体，覆盖 2010—2026\\全部英语二 Part A 真题}
\end{minipage}\hfill
\begin{minipage}[t]{0.32\linewidth}
{\color{accent}\fontsize{30}{32}\selectfont\bfseries 95–105}\par
{\small\color{gray} 词，每篇正文都按考场\\词数控制，逐篇统计}
\end{minipage}\hfill
\begin{minipage}[t]{0.32\linewidth}
{\color{accent}\fontsize{30}{32}\selectfont\bfseries 5}\par
{\small\color{gray} 个空，考场只改活动、\\时间、地点、理由、要求}
\end{minipage}\par

\vspace{22pt}
\noindent\colorbox{soft}{\parbox{\dimexpr\linewidth-2\fboxsep\relax}{%%
\vspace{5pt}
{\small\bfseries 这本书怎么用}\par\vspace{3pt}
{\footnotesize
一、先看第 2 页：分从哪来、格式怎么摆、三段各写几句。\par
二、再按第 3 页的真题地图，找到你要练的那一篇。\par
三、每篇只背“整篇模板”，然后用同一页的真题落地对照，看五个空怎么填。\par
四、写完用第 16 页的 60 秒自评表打勾，再按批改标准给自己定档。\par}
\vspace{5pt}}}

\vspace{12pt}
\noindent{\footnotesize\color{gray}
本书与市面“万能模板”的区别：不给你一堆可替换句型，而是给你 12 篇完整信件；\\每篇的骨架、衔接、语气是固定的，考场上只换信息，不现编句子。\\[3pt]
全部英文为原创；骨架按文体逐篇重写，不照搬南山《小作文万用模版》的万能句。}

\vspace{18pt}
{\color{accent}\bfseries 12 篇目录}\par\vspace{2pt}
{\color{rule}\rule{\linewidth}{0.5pt}}\par\vspace{3pt}
{\footnotesize
\noindent\begin{tabularx}{\linewidth}{@{}XXX@{}}
01\quad 邀请信 & 02\quad 接受邀请 & 03\quad 婉拒邀请 \\
04\quad 建议信 & 05\quad 感谢信 & 06\quad 祝贺信 \\
07\quad 道歉信 & 08\quad 投诉信 & 09\quad 咨询信 \\
10\quad 通知 & 11\quad 介绍·告知 & 12\quad 推荐信 \\
\end{tabularx}}\par
\vspace{2pt}
{\footnotesize\color{gray} 页码见页脚；先练 01、04、11 三篇（近五年高频），再补其余。}\par

\vspace{20pt}
{\color{rule}\rule{\linewidth}{0.5pt}}\par\vspace{8pt}
\noindent\begin{minipage}[t]{0.235\linewidth}
{\color{accent}\bfseries ① 定文体}\par{\scriptsize\color{gray} 第 3 页真题地图\\找动作词}
\end{minipage}\hfill
\begin{minipage}[t]{0.235\linewidth}
{\color{accent}\bfseries ② 背整篇}\par{\scriptsize\color{gray} 每页第一节\\完整信件}
\end{minipage}\hfill
\begin{minipage}[t]{0.235\linewidth}
{\color{accent}\bfseries ③ 换五词}\par{\scriptsize\color{gray} 同页换装表\\只改槽位}
\end{minipage}\hfill
\begin{minipage}[t]{0.235\linewidth}
{\color{accent}\bfseries ④ 自评}\par{\scriptsize\color{gray} 第 16 页\\60 秒清单}
\end{minipage}\par

\vfill
{\footnotesize\color{gray} 编写依据：考研英语（二）考试大纲对 Part A 的要求；评分档位按工作区《批改指令》的民间参考标准整理，不人为封顶。}\par


\clearpage
\vspace*{-4pt}
{\color{accent}\bfseries\LARGE 考场先看这一页}\par\vspace{3pt}
{\color{accent}\rule{\linewidth}{1.1pt}}\par
\vspace{5pt}
{\small\color{gray} 小作文 10 分，15 分钟。先保证“任务完成＋格式正确”，再谈语言漂亮。}\par
\vspace{8pt}

{\color{accent}\bfseries 一、阅卷人按四个维度给分}\par
{\color{rule}\rule{\linewidth}{0.5pt}}\par\vspace{2pt}
{\footnotesize
\begin{tabularx}{\linewidth}{@{}>{\raggedright\arraybackslash}p{2.0cm}>{\raggedright\arraybackslash}p{5.4cm}X@{}}
\toprule
{\color{accent}\bfseries 维度} & {\color{accent}\bfseries 卷面上看什么} & {\color{accent}\bfseries 你写的时候怎么做} \\
\midrule
任务完成 & 题干两个动作都有落点 & 动笔前把 1) 2) 抄在草稿纸，写完逐条打勾 \\
语言准确 & 主谓一致、时态、搭配、拼写 & 只用有把握的句型；不确定的词换简单词 \\
格式语气 & 称呼、结束语、署名、通知标题落款 & 按下面两张表套，30 秒能查完 \\
连贯变化 & 段间有过渡，句式不重复 & 每段首句换开头；I 开头不超过四句 \\
\bottomrule
\end{tabularx}}\par
\vspace{3pt}
{\footnotesize\color{gray} 四维里“任务完成”权重最高：漏一个动作，语言再好也上不了 7 分。}\par

\vspace{9pt}
{\color{accent}\bfseries 二、书信格式四要素（30 秒能查完）}\par
{\color{rule}\rule{\linewidth}{0.5pt}}\par\vspace{2pt}
{\footnotesize
\begin{tabularx}{\linewidth}{@{}>{\raggedright\arraybackslash}p{1.8cm}>{\raggedright\arraybackslash}p{6.2cm}X@{}}
\toprule
{\color{accent}\bfseries 位置} & {\color{accent}\bfseries 写法} & {\color{accent}\bfseries 例} \\
\midrule
称呼 & 顶格，逗号结尾；熟人用名，师长用姓 & Dear Jack, / Dear Professor Wang, \\
正文 & 2—3 段，每段一个功能，段首不空格 & 见各文体页 \\
结束语 & 顶格，逗号结尾，与署名同一语气层 & Yours sincerely, / Best wishes, \\
署名 & 另起一行，写题目指定姓名 & Li Ming / Zhang Wei \\
\bottomrule
\end{tabularx}}\par

\vspace{9pt}
{\color{accent}\bfseries 三、通知的三处不同}\par
{\color{rule}\rule{\linewidth}{0.5pt}}\par\vspace{2pt}
{\footnotesize
\begin{tabularx}{\linewidth}{@{}>{\raggedright\arraybackslash}p{1.8cm}>{\raggedright\arraybackslash}p{6.2cm}X@{}}
\toprule
{\color{accent}\bfseries 位置} & {\color{accent}\bfseries 写法} & {\color{accent}\bfseries 例} \\
\midrule
标题 & 居中，单独一行 & Notice \\
称呼 & 不写称呼 & — \\
落款 & 右下角两行：单位＋日期 & The Student Union / June 20, 2026 \\
\bottomrule
\end{tabularx}}\par

\vspace{9pt}
{\color{accent}\bfseries 四、三段各写几句}\par
{\color{rule}\rule{\linewidth}{0.5pt}}\par\vspace{2pt}
{\footnotesize
\begin{itemize}[leftmargin=1.2em,itemsep=1pt,topsep=1pt,parsep=0pt]
\item {\bfseries 第一段（1—2 句）}：写作目的＋事由。不写“I am a student”，不铺背景。
\item {\bfseries 第二段（3—4 句）}：全部任务要点，一个要点一句。合并题用 As for… / As to… 切开。
\item {\bfseries 第三段（1—2 句）}：期待、感谢、请回复，三选一，不重复第二段。
\end{itemize}}\par

\vspace{7pt}
{\color{accent}\bfseries 五、词数与时间}\par
{\color{rule}\rule{\linewidth}{0.5pt}}\par\vspace{2pt}
{\footnotesize
\begin{itemize}[leftmargin=1.2em,itemsep=1pt,topsep=1pt,parsep=0pt]
\item 正文 95—105 词，7—9 句；{\bfseries 称呼、结束语、署名不算}，也不写地址和日期（通知除外）。
\item 每行 10—12 词，卷面 8—10 行；超过 120 词不加分，只放大错误。
\item 15 分钟：审题画要点 3 分钟 → 写 10 分钟 → 检查 2 分钟。
\end{itemize}}\par

\vspace{7pt}
{\color{accent}\bfseries 六、六条铁律}\par
{\color{rule}\rule{\linewidth}{0.5pt}}\par\vspace{2pt}
{\footnotesize
\begin{enumerate}[leftmargin=1.4em,itemsep=0.5pt,topsep=1pt,parsep=0pt]
\item 落款用题目给的姓名（Li Ming / Zhang Wei），不写自己的名字。
\item 题干给了时间、地点、姓名，一律照抄；没给就补一个具体而不夸张的（next Saturday / Room 203）。
\item 合并题必须两段分明，段首用 As for… / As to… 提示阅卷人“第二个动作在这里”。
\item 不确定的难词一律换简单词；拼写错误比用词普通更扣分。
\item 称呼：熟人 Dear Jack,；师长 Dear Professor Wang,；不知姓名 Dear Sir or Madam,。
\item 通知不写称呼，标题居中，落款单位＋日期在右下角。
\end{enumerate}}\par


\clearpage
\vspace*{-4pt}
{\color{accent}\bfseries\LARGE 真题地图：先定位，再动笔}\par\vspace{3pt}
{\color{accent}\rule{\linewidth}{1.1pt}}\par
\vspace{5pt}
{\small\color{gray} 题干永远给两件事：写什么（事务）＋要你完成什么动作。先找动作词，文体就定了。}\par
\vspace{5pt}
{\scriptsize
\begin{longtable}{@{}>{\raggedright\arraybackslash}p{1.15cm}>{\raggedright\arraybackslash}p{3.05cm}>{\raggedright\arraybackslash}p{7.9cm}>{\raggedright\arraybackslash}p{2.15cm}@{}}
\toprule
{\color{accent}\bfseries 年份} & {\color{accent}\bfseries 文体} & {\color{accent}\bfseries 题干两个动作} & {\color{accent}\bfseries 用哪篇} \\
\midrule
\endfirsthead
\toprule
{\color{accent}\bfseries 年份} & {\color{accent}\bfseries 文体} & {\color{accent}\bfseries 题干两个动作} & {\color{accent}\bfseries 用哪篇} \\
\midrule
\endhead
2010 & 感谢＋邀请 & 感谢接待；欢迎访华 & 05 \\
2011 & 祝贺＋建议 & 祝贺录取；适应大学生活的建议 & 06 \\
2012 & 投诉 & 电子词典故障；要求尽快解决 & 08 \\
2013 & 通知 & 慈善义卖细节；鼓励同学参加 & 10 \\
2014 & 咨询＋告知 & 介绍生活习惯；询问当地生活 & 09 \\
2015 & 通知 & 夏令营活动；招募志愿者 & 10 \\
2016 & 回复祝贺＋建议 & 感谢祝贺；给出翻译学习建议 & 06 · 04 \\
2017 & 回复邀请·接受 & 接受邀请；介绍讲座要点 & 02 \\
2018 & 道歉＋另约 & 取消拜访；解释并另约 & 07 \\
2019 & 建议＋安排 & 辩论题目及理由；告知安排 & 04 \\
2020 & 介绍＋提示 & 介绍历史景点；给出参观提示 & 11 \\
2021 & 邀请 & 邀请参加线上会议；告知细节 & 01 \\
2022 & 介绍＋邀请 & 介绍校园美食节；邀请参加 & 11 · 01 \\
2023 & 建议＋理由 & 艺术展与机器人展二选一；给理由 & 04 \\
2024 & 建议＋征询 & 提出调查计划；征求对方意见 & 04 · 11 \\
2025 & 介绍＋邀请 & 介绍短剧；邀请朋友参加 & 11 · 01 \\
2026* & 感谢＋询问 & 谈对旅行视频的感受；想了解更多 & 05 · 09 \\
\bottomrule
\end{longtable}}\par
\vspace{4pt}
{\footnotesize\color{gray} * 2026 为回忆版，以当年正式试题为准。}\par

\vspace{12pt}
{\color{accent}\bfseries 三条读表结论}\par
{\color{rule}\rule{\linewidth}{0.5pt}}\par\vspace{4pt}
{\footnotesize
\begin{enumerate}[leftmargin=1.4em,rightmargin=6pt,itemsep=2pt,topsep=1pt,parsep=0pt]
\item {\bfseries 十二类足够覆盖全部真题。}反复考的只有邀请、回复、建议、感谢、祝贺、道歉、投诉、咨询、通知、介绍；推荐与婉拒备作模拟。
\item {\bfseries 一半年份是“合并题”。}2010、2011、2014、2016、2019、2024 都是两个动作；写的时候第二段开头用 As for… 切开，阅卷人找得到第二个要点。
\item {\bfseries 近五年偏爱“介绍＋邀请 / 建议＋理由”。}2021—2025 全部落在 01、04、11 三篇里；这三篇要练到不看模板能写。
\end{enumerate}}\par

\vspace{8pt}
{\color{accent}\bfseries 合并题的五个现成切法}\par
{\color{rule}\rule{\linewidth}{0.5pt}}\par\vspace{4pt}
{\footnotesize
\begin{itemize}[leftmargin=1.2em,itemsep=2pt,topsep=1pt,parsep=0pt]
\item {\bfseries 2010 感谢＋邀请}：第二段接 05 的 By the way, you are welcome to…
\item {\bfseries 2011 祝贺＋建议}：第二段接 06 的 As you are about to… I would like to offer two suggestions.
\item {\bfseries 2014 告知＋咨询}：第二段接 09 的 To give you some idea of my habits…
\item {\bfseries 2019 建议＋安排}：第二段另起 As for the arrangements, the debate will take place…
\item {\bfseries 2024 计划＋征询}：结尾加一句 I wonder what you think of it.
\end{itemize}}\par

\vspace{8pt}
{\color{accent}\bfseries 文体频次（2010—2026，含合并题）}\par
{\color{rule}\rule{\linewidth}{0.5pt}}\par\vspace{4pt}
{\footnotesize\renewcommand{\arraystretch}{1.0}
\begin{tabularx}{\linewidth}{@{}>{\raggedright\arraybackslash}p{2.6cm}>{\raggedright\arraybackslash}p{0.9cm}>{\raggedright\arraybackslash}p{0.9cm}>{\raggedright\arraybackslash}p{6.4cm}X@{}}
\toprule
{\color{accent}\bfseries 文体} & {\color{accent}\bfseries 次数} & {\color{accent}\bfseries 篇} & {\color{accent}\bfseries 年份} & {\color{accent}\bfseries 频次条} \\
\midrule
建议 & 5 & 04 & 2011·2016·2019·2023·2024 & {\color{accent}■■■■■} \\
介绍·告知 & 5 & 11 & 2014·2020·2022·2024·2025 & {\color{accent}■■■■■} \\
邀请 & 4 & 01 & 2010·2021·2022·2025 & {\color{accent}■■■■} \\
感谢 & 3 & 05 & 2010·2016·2026* & {\color{accent}■■■} \\
通知 & 2 & 10 & 2013·2015 & {\color{accent}■■} \\
祝贺 & 2 & 06 & 2011·2016 & {\color{accent}■■} \\
咨询 & 2 & 09 & 2014·2026* & {\color{accent}■■} \\
投诉 & 1 & 08 & 2012 & {\color{accent}■} \\
道歉 & 1 & 07 & 2018 & {\color{accent}■} \\
\bottomrule
\end{tabularx}}\par


\clearpage
\phantomsection\label{g01}
\noindent{\color{accent}\rule{\linewidth}{1.5pt}}\par\vspace{3pt}
\noindent{\fontsize{26}{28}\selectfont\bfseries\color{accent}01}\hspace{6pt}
{\LARGE\bfseries 邀请信}\hspace{7pt}{\large\color{gray} Invitation}\hfill
{\footnotesize\color{gray} 真题 2021 · 2022 · 2025}\par
\vspace{1pt}
{\color{accent!35}\rule{\linewidth}{0.5pt}}\par\vspace{3pt}
{\footnotesize {\color{accent}\bfseries 题干信号}\quad 题干出现 invite / 邀请……参加，并要求交代时间、地点或安排}\par
{\footnotesize {\color{accent}\bfseries 写作身份}\quad 活动组织者、班委或社团负责人；落款写题目指定姓名}\par
{\footnotesize {\color{accent}\bfseries 必须写到}\quad ①活动名称与目的　②时间与地点　③两个环节　④请对方在某日前回复}\par
{\footnotesize {\color{accent}\bfseries 一句话策略}\quad 第一段一句点题，第二段全部交代安排，第三段只做一件事：促回复}\par

\vspace{5pt}
{\color{accent}\bfseries 一、整篇模板（背这一篇）}\hfill{\footnotesize\color{gray} 正文 98 词 · 7 句}\par
\vspace{2pt}
\noindent\fcolorbox{rule}{soft}{\begin{minipage}{\dimexpr\linewidth-2\fboxsep-2\fboxrule\relax}\vspace{3pt}\small\linespread{1.03}\selectfont Dear {\color{hot}\bfseries 【称呼】},\par
\vspace{3pt}
{\color{accent}\bfseries ①}I am writing to invite you to {\color{hot}\bfseries 【活动名称】}, which {\color{hot}\bfseries 【主办方】} is organizing for {\color{hot}\bfseries 【参加对象】}.\par
\vspace{2pt}
{\color{accent}\bfseries ②}The event will take place {\color{hot}\bfseries 【日期】} at {\color{hot}\bfseries 【地点】}.{\color{accent}\bfseries ③}It will open with {\color{hot}\bfseries 【环节一】}, followed by {\color{hot}\bfseries 【环节二】}, and end with {\color{hot}\bfseries 【环节三】}.{\color{accent}\bfseries ④}There is no charge, and {\color{hot}\bfseries 【附加安排】}.{\color{accent}\bfseries ⑤}You will also have a chance to {\color{hot}\bfseries 【参与收获】}.\par
\vspace{2pt}
{\color{accent}\bfseries ⑥}I do hope you can join us.{\color{accent}\bfseries ⑦}Would you please let me know by {\color{hot}\bfseries 【回复期限】} whether you are free that day?\par
\vspace{2pt}
\vspace{2pt}Yours sincerely,\par
Li Ming\par\vspace{3pt}\end{minipage}}\par
\vspace{1pt}
{\footnotesize\color{gray} 段落分工｜① 点题｜②—⑤ 时间、地点、环节、收获｜⑥—⑦ 促回复}\par

\vspace{6pt}
{\color{accent}\bfseries 二、真题落地（2022 真题 · 校园美食节）}\par
\vspace{2pt}
\noindent\fcolorbox{sandline}{sand}{\begin{minipage}{\dimexpr\linewidth-2\fboxsep-2\fboxrule\relax}\vspace{3pt}\small\linespread{1.03}\selectfont Dear Jack,\par
\vspace{3pt}
{\color{accent}\bfseries ①}I am writing to invite you to our Campus Food Festival, which the Student Union is organizing for international students.\par
\vspace{2pt}
{\color{accent}\bfseries ②}The event will take place next Saturday at the Central Square. {\color{accent}\bfseries ③}It will open with a short cooking show, followed by a tasting session, and end with a vote for the most popular dish. {\color{accent}\bfseries ④}There is no charge, and light refreshments will be served. {\color{accent}\bfseries ⑤}You will also have a chance to try dishes from over ten provinces.\par
\vspace{2pt}
{\color{accent}\bfseries ⑥}I do hope you can join us. {\color{accent}\bfseries ⑦}Would you please let me know by Wednesday whether you are free that day?\par
\vspace{2pt}
\vspace{2pt}Yours sincerely,\par
Li Ming\par\vspace{3pt}\end{minipage}}\par

\vspace{6pt}
{\color{accent}\bfseries 三、换装表（只改这几处）}\par
\vspace{2pt}
{\footnotesize
\begin{tabularx}{\linewidth}{@{}>{\raggedright\arraybackslash}p{3.5cm}>{\raggedright\arraybackslash}p{5.6cm}X@{}}
\toprule
{\color{accent}\bfseries 槽位} & {\color{accent}\bfseries 本篇真题怎么填} & {\color{accent}\bfseries 可换词块} \\
\midrule
【活动名称】 & our Campus Food Festival & a charity sale / a welcome party \\
【日期】＋【地点】 & next Saturday / the Central Square & next Friday at 7 p.m. / Room 203 \\
【环节一】【环节二】【环节三】 & cooking show → tasting → vote & opening speech → game → free talk \\
【参与收获】 & try dishes from over ten provinces & meet students from other majors \\
\bottomrule
\end{tabularx}}\par

\vspace{5pt}
\noindent\begin{minipage}[t]{0.555\linewidth}
{\color{accent}\bfseries 四、阅卷人找什么}\par
{\footnotesize\begin{itemize}[leftmargin=1.2em,itemsep=1pt,topsep=1pt,parsep=0pt]\setlength{\leftmargin}{1.1em}
\item 活动名称、时间、地点三样齐全，且都在第二段，阅卷人一眼能找到
\item 两个以上具体环节（不是“interesting activities”这种空话）
\item 结尾有一句明确的回复请求，并给出期限
\end{itemize}}
\end{minipage}\hfill
\begin{minipage}[t]{0.415\linewidth}
{\color{hot}\bfseries 别踩这两个坑}\par
{\footnotesize\begin{itemize}[leftmargin=1.2em,itemsep=1pt,topsep=1pt,parsep=0pt]\setlength{\leftmargin}{1.1em}
\item 只写“下周六”，不写几点、在哪 —— 直接算要点不全
\item 通篇 I will… I will…，把邀请信写成日程表
\end{itemize}}
\end{minipage}\par

\vspace{6pt}
{\color{accent}\bfseries 五、想上 9 分，从这两句里挑一句加在结尾}\par
\vspace{2pt}
{\footnotesize
\begin{itemize}[leftmargin=1.2em,itemsep=1.5pt,topsep=1pt,parsep=0pt]
\item Should you need any further details, please feel free to contact me.\quad{\color{gray}（条件倒装，替代 if you need…，结尾更像公务信）}
\item I would be more than happy to reserve a seat for you.\quad{\color{gray}（would be more than happy 比 I am happy 礼貌一档）}
\end{itemize}}\par

\vspace{6pt}
{\color{accent}\bfseries 六、题干变体怎么改（同一篇，换两个词）}\par
\vspace{2pt}
{\footnotesize
\begin{itemize}[leftmargin=1.2em,itemsep=1.5pt,topsep=1pt,parsep=0pt]
\item 2021 线上会议：地点改 on Zoom，环节改 agenda items，附上会议号。
\item 2025 短剧：主办方改 the Drama Club，参加对象改 your classmates。
\end{itemize}}\par


\clearpage
\phantomsection\label{g02}
\noindent{\color{accent}\rule{\linewidth}{1.5pt}}\par\vspace{3pt}
\noindent{\fontsize{26}{28}\selectfont\bfseries\color{accent}02}\hspace{6pt}
{\LARGE\bfseries 回复邀请 · 接受}\hspace{7pt}{\large\color{gray} Accepting an Invitation}\hfill
{\footnotesize\color{gray} 真题 2017（回复类真题）}\par
\vspace{1pt}
{\color{accent!35}\rule{\linewidth}{0.5pt}}\par\vspace{3pt}
{\footnotesize {\color{accent}\bfseries 题干信号}\quad 题干出现 reply to / 回复某封邀请信，第一个动作是接受}\par
{\footnotesize {\color{accent}\bfseries 写作身份}\quad 受邀者；称呼沿用邀请人的身份与姓氏}\par
{\footnotesize {\color{accent}\bfseries 必须写到}\quad ①感谢邀请　②明确接受　③讲座/发言要点　④到场准备或需对方配合的事}\par
{\footnotesize {\color{accent}\bfseries 一句话策略}\quad 第一段必须出现 accept 这个词；第二段用“要点＋准备”证明你真的要来}\par

\vspace{5pt}
{\color{accent}\bfseries 一、整篇模板（背这一篇）}\hfill{\footnotesize\color{gray} 正文 98 词 · 7 句}\par
\vspace{2pt}
\noindent\fcolorbox{rule}{soft}{\begin{minipage}{\dimexpr\linewidth-2\fboxsep-2\fboxrule\relax}\vspace{3pt}\small\linespread{1.03}\selectfont Dear {\color{hot}\bfseries 【称呼】},\par
\vspace{3pt}
{\color{accent}\bfseries ①}Thank you for your invitation to {\color{hot}\bfseries 【活动】}.{\color{accent}\bfseries ②}I am glad to accept it and look forward to {\color{hot}\bfseries 【期待】}.\par
\vspace{2pt}
{\color{accent}\bfseries ③}As you suggested, I will center my talk on {\color{hot}\bfseries 【要点一】} and {\color{hot}\bfseries 【要点二】}.{\color{accent}\bfseries ④}I plan to arrive at {\color{hot}\bfseries 【地点】} half an hour early so that I can {\color{hot}\bfseries 【准备工作】}.{\color{accent}\bfseries ⑤}If you need any materials from me, please tell me before {\color{hot}\bfseries 【日期】}.\par
\vspace{2pt}
{\color{accent}\bfseries ⑥}Thank you again for thinking of me.{\color{accent}\bfseries ⑦}I will do my best to make the talk useful to {\color{hot}\bfseries 【听众】}, and I will bring a short handout.\par
\vspace{2pt}
\vspace{2pt}Yours sincerely,\par
Li Ming\par\vspace{3pt}\end{minipage}}\par
\vspace{1pt}
{\footnotesize\color{gray} 段落分工｜①—② 感谢＋接受｜③—⑤ 要点与准备｜⑥—⑦ 再谢＋承诺}\par

\vspace{6pt}
{\color{accent}\bfseries 二、真题落地（2017 真题 · 接受邀请并介绍讲座要点）}\par
\vspace{2pt}
\noindent\fcolorbox{sandline}{sand}{\begin{minipage}{\dimexpr\linewidth-2\fboxsep-2\fboxrule\relax}\vspace{3pt}\small\linespread{1.03}\selectfont Dear Professor Williams,\par
\vspace{3pt}
{\color{accent}\bfseries ①}Thank you for your invitation to give a presentation on Chinese culture. {\color{accent}\bfseries ②}I am glad to accept it and look forward to meeting your students.\par
\vspace{2pt}
{\color{accent}\bfseries ③}As you suggested, I will center my talk on traditional festivals and table manners. {\color{accent}\bfseries ④}I plan to arrive at the lecture hall half an hour early so that I can test the projector. {\color{accent}\bfseries ⑤}If you need any materials from me, please tell me before Friday.\par
\vspace{2pt}
{\color{accent}\bfseries ⑥}Thank you again for thinking of me. {\color{accent}\bfseries ⑦}I will do my best to make the talk useful to the international students, and I will bring a short handout.\par
\vspace{2pt}
\vspace{2pt}Yours sincerely,\par
Li Ming\par\vspace{3pt}\end{minipage}}\par

\vspace{6pt}
{\color{accent}\bfseries 三、换装表（只改这几处）}\par
\vspace{2pt}
{\footnotesize
\begin{tabularx}{\linewidth}{@{}>{\raggedright\arraybackslash}p{3.5cm}>{\raggedright\arraybackslash}p{5.6cm}X@{}}
\toprule
{\color{accent}\bfseries 槽位} & {\color{accent}\bfseries 本篇真题怎么填} & {\color{accent}\bfseries 可换词块} \\
\midrule
【活动】 & give a presentation on Chinese culture & attend the seminar / join the panel \\
【要点一】【要点二】 & traditional festivals / table manners & food culture / family life \\
【准备工作】 & test the projector & check the microphone / print handouts \\
【听众】 & the international students & first-year students / the audience \\
\bottomrule
\end{tabularx}}\par

\vspace{5pt}
\noindent\begin{minipage}[t]{0.555\linewidth}
{\color{accent}\bfseries 四、阅卷人找什么}\par
{\footnotesize\begin{itemize}[leftmargin=1.2em,itemsep=1pt,topsep=1pt,parsep=0pt]\setlength{\leftmargin}{1.1em}
\item 第一段有明确态度词（accept / delighted to accept），不拖到第二段
\item 讲座要点具体到两个名词，不是“something about culture”
\item 主动交代到场时间与准备事项，体现“真接受”
\end{itemize}}
\end{minipage}\hfill
\begin{minipage}[t]{0.415\linewidth}
{\color{hot}\bfseries 别踩这两个坑}\par
{\footnotesize\begin{itemize}[leftmargin=1.2em,itemsep=1pt,topsep=1pt,parsep=0pt]\setlength{\leftmargin}{1.1em}
\item 只写 I am glad to accept，后面全是空话，第二个动作没落地
\item 把接受信写成感谢信，通篇 thank you，没有安排
\end{itemize}}
\end{minipage}\par

\vspace{6pt}
{\color{accent}\bfseries 五、想上 9 分，从这两句里挑一句加在结尾}\par
\vspace{2pt}
{\footnotesize
\begin{itemize}[leftmargin=1.2em,itemsep=1.5pt,topsep=1pt,parsep=0pt]
\item I have already prepared a short handout for the audience.\quad{\color{gray}（现在完成时，证明“我真的准备好了”）}
\item Please let me know if there is anything else I should bring.\quad{\color{gray}（主动询问，显得可靠，不像背模板）}
\end{itemize}}\par

\vspace{6pt}
{\color{accent}\bfseries 六、题干变体怎么改（同一篇，换两个词）}\par
\vspace{2pt}
{\footnotesize
\begin{itemize}[leftmargin=1.2em,itemsep=1.5pt,topsep=1pt,parsep=0pt]
\item 线上会议变体：准备工作改 test my microphone，日期改 the day before。
\item 要点部分改成两个具体话题，不要写 many interesting things。
\end{itemize}}\par


\clearpage
\phantomsection\label{g03}
\noindent{\color{accent}\rule{\linewidth}{1.5pt}}\par\vspace{3pt}
\noindent{\fontsize{26}{28}\selectfont\bfseries\color{accent}03}\hspace{6pt}
{\LARGE\bfseries 回复邀请 · 婉拒}\hspace{7pt}{\large\color{gray} Declining an Invitation}\hfill
{\footnotesize\color{gray} 真题 模拟题 · 大纲内文体}\par
\vspace{1pt}
{\color{accent!35}\rule{\linewidth}{0.5pt}}\par\vspace{3pt}
{\footnotesize {\color{accent}\bfseries 题干信号}\quad 题干要求 decline / 不能参加，并说明原因或另约}\par
{\footnotesize {\color{accent}\bfseries 写作身份}\quad 受邀者；语气要礼貌，但不能含糊}\par
{\footnotesize {\color{accent}\bfseries 必须写到}\quad ①感谢邀请　②明确不能参加　③一句真实原因　④替代时间或替代方案}\par
{\footnotesize {\color{accent}\bfseries 一句话策略}\quad 先谢、再拒、给原因、给替代——四步缺一不可，尤其不能只拒绝}\par

\vspace{5pt}
{\color{accent}\bfseries 一、整篇模板（背这一篇）}\hfill{\footnotesize\color{gray} 正文 100 词 · 8 句}\par
\vspace{2pt}
\noindent\fcolorbox{rule}{soft}{\begin{minipage}{\dimexpr\linewidth-2\fboxsep-2\fboxrule\relax}\vspace{3pt}\small\linespread{1.03}\selectfont Dear {\color{hot}\bfseries 【称呼】},\par
\vspace{3pt}
{\color{accent}\bfseries ①}Thank you for inviting me to {\color{hot}\bfseries 【活动】}.{\color{accent}\bfseries ②}I am sorry to say that I shall not be able to attend it, though I really wish I could.\par
\vspace{2pt}
{\color{accent}\bfseries ③}The problem is that {\color{hot}\bfseries 【冲突原因】} on that very day, and I cannot change the arrangement.{\color{accent}\bfseries ④}I feel bad about missing {\color{hot}\bfseries 【期待内容】}, which I have been looking forward to.{\color{accent}\bfseries ⑤}Could we meet at another time?{\color{accent}\bfseries ⑥}I would be free {\color{hot}\bfseries 【备选时间】}, and I am happy to {\color{hot}\bfseries 【替代方案】}.\par
\vspace{2pt}
{\color{accent}\bfseries ⑦}I hope you will understand.{\color{accent}\bfseries ⑧}Thank you again for your kindness.\par
\vspace{2pt}
\vspace{2pt}Yours sincerely,\par
Li Ming\par\vspace{3pt}\end{minipage}}\par
\vspace{1pt}
{\footnotesize\color{gray} 段落分工｜①—② 感谢＋拒绝｜③—⑥ 原因＋替代｜⑦—⑧ 致歉收尾}\par

\vspace{6pt}
{\color{accent}\bfseries 二、真题落地（模拟题 · 与考试冲突，婉拒研讨会）}\par
\vspace{2pt}
\noindent\fcolorbox{sandline}{sand}{\begin{minipage}{\dimexpr\linewidth-2\fboxsep-2\fboxrule\relax}\vspace{3pt}\small\linespread{1.03}\selectfont Dear Professor Wang,\par
\vspace{3pt}
{\color{accent}\bfseries ①}Thank you for inviting me to the seminar on urban planning. {\color{accent}\bfseries ②}I am sorry to say that I shall not be able to attend it, though I really wish I could.\par
\vspace{2pt}
{\color{accent}\bfseries ③}The problem is that I have to sit for an important exam on that very day, and I cannot change the arrangement. {\color{accent}\bfseries ④}I feel bad about missing your talk, which I have been looking forward to. {\color{accent}\bfseries ⑤}Could we meet at another time? {\color{accent}\bfseries ⑥}I would be free next Friday afternoon, and I am happy to visit your office then.\par
\vspace{2pt}
{\color{accent}\bfseries ⑦}I hope you will understand. {\color{accent}\bfseries ⑧}Thank you again for your kindness.\par
\vspace{2pt}
\vspace{2pt}Yours sincerely,\par
Li Ming\par\vspace{3pt}\end{minipage}}\par

\vspace{6pt}
{\color{accent}\bfseries 三、换装表（只改这几处）}\par
\vspace{2pt}
{\footnotesize
\begin{tabularx}{\linewidth}{@{}>{\raggedright\arraybackslash}p{3.5cm}>{\raggedright\arraybackslash}p{5.6cm}X@{}}
\toprule
{\color{accent}\bfseries 槽位} & {\color{accent}\bfseries 本篇真题怎么填} & {\color{accent}\bfseries 可换词块} \\
\midrule
【活动】 & the seminar on urban planning & the conference / the interview \\
【冲突原因】 & I have to sit for an important exam & I will be away on a field trip \\
【替代方案】 & visit your office then & read your paper and send comments \\
\bottomrule
\end{tabularx}}\par

\vspace{5pt}
\noindent\begin{minipage}[t]{0.555\linewidth}
{\color{accent}\bfseries 四、阅卷人找什么}\par
{\footnotesize\begin{itemize}[leftmargin=1.2em,itemsep=1pt,topsep=1pt,parsep=0pt]\setlength{\leftmargin}{1.1em}
\item 拒绝得明确（shall not be able to attend），不是“可能不行”
\item 原因只有一句，不展开诉苦；替代方案必须具体到时间
\item 语气礼貌但不卑微，不用 sorry 开头连写三句
\end{itemize}}
\end{minipage}\hfill
\begin{minipage}[t]{0.415\linewidth}
{\color{hot}\bfseries 别踩这两个坑}\par
{\footnotesize\begin{itemize}[leftmargin=1.2em,itemsep=1pt,topsep=1pt,parsep=0pt]\setlength{\leftmargin}{1.1em}
\item 只拒绝、不给替代方案 —— 任务只完成一半
\item 原因写成一段故事（家里、身体、天气全上），挤掉替代安排
\end{itemize}}
\end{minipage}\par

\vspace{6pt}
{\color{accent}\bfseries 五、想上 9 分，从这两句里挑一句加在结尾}\par
\vspace{2pt}
{\footnotesize
\begin{itemize}[leftmargin=1.2em,itemsep=1.5pt,topsep=1pt,parsep=0pt]
\item I would be grateful if we could fix another time that suits you.\quad{\color{gray}（礼貌请求，替代生硬的 Can we meet again?）}
\item I hope this will not cause you too much inconvenience.\quad{\color{gray}（先顾对方感受，婉拒信的关键一句）}
\end{itemize}}\par

\vspace{6pt}
{\color{accent}\bfseries 六、题干变体怎么改（同一篇，换两个词）}\par
\vspace{2pt}
{\footnotesize
\begin{itemize}[leftmargin=1.2em,itemsep=1.5pt,topsep=1pt,parsep=0pt]
\item 只用于“被邀请但去不了”；2018 那种“自己取消约定”用 07 道歉信。
\item 婉拒教授时保留 Professor 姓氏，不要用 Dear teacher。
\end{itemize}}\par


\clearpage
\phantomsection\label{g04}
\noindent{\color{accent}\rule{\linewidth}{1.5pt}}\par\vspace{3pt}
\noindent{\fontsize{26}{28}\selectfont\bfseries\color{accent}04}\hspace{6pt}
{\LARGE\bfseries 建议信}\hspace{7pt}{\large\color{gray} Advice / Suggestion}\hfill
{\footnotesize\color{gray} 真题 2011 · 2016 · 2019 · 2023 · 2024}\par
\vspace{1pt}
{\color{accent!35}\rule{\linewidth}{0.5pt}}\par\vspace{3pt}
{\footnotesize {\color{accent}\bfseries 题干信号}\quad 题干出现 advice / suggestions / which one / 征求意见，要你给主张}\par
{\footnotesize {\color{accent}\bfseries 写作身份}\quad 朋友、同学或下属；同级用商量口气，对师长用 I would suggest}\par
{\footnotesize {\color{accent}\bfseries 必须写到}\quad ①明确建议哪一个　②两条理由　③一条附加好处　④把决定权还给对方}\par
{\footnotesize {\color{accent}\bfseries 一句话策略}\quad 第一段一句话给结论，第二段两条理由各一句，第三段客气收尾}\par

\vspace{5pt}
{\color{accent}\bfseries 一、整篇模板（背这一篇）}\hfill{\footnotesize\color{gray} 正文 99 词 · 8 句}\par
\vspace{2pt}
\noindent\fcolorbox{rule}{soft}{\begin{minipage}{\dimexpr\linewidth-2\fboxsep-2\fboxrule\relax}\vspace{3pt}\small\linespread{1.03}\selectfont Dear {\color{hot}\bfseries 【称呼】},\par
\vspace{3pt}
{\color{accent}\bfseries ①}You asked me for my advice on {\color{hot}\bfseries 【咨询事项】}.{\color{accent}\bfseries ②}I would suggest {\color{hot}\bfseries 【建议对象】}.\par
\vspace{2pt}
{\color{accent}\bfseries ③}For one thing, {\color{hot}\bfseries 【理由一】}.{\color{accent}\bfseries ④}For another, {\color{hot}\bfseries 【理由二】}.{\color{accent}\bfseries ⑤}Besides, {\color{hot}\bfseries 【附加好处】}, which {\color{hot}\bfseries 【对方获益】}.\par
\vspace{2pt}
{\color{accent}\bfseries ⑥}Of course, the final choice is yours.{\color{accent}\bfseries ⑦}I have seen the poster, and it looks really attractive.{\color{accent}\bfseries ⑧}If you decide to go, I would be happy to {\color{hot}\bfseries 【陪同安排】}, and I am sure you will enjoy it.\par
\vspace{2pt}
\vspace{2pt}Yours sincerely,\par
Li Ming\par\vspace{3pt}\end{minipage}}\par
\vspace{1pt}
{\footnotesize\color{gray} 段落分工｜①—② 给结论｜③—⑤ 两条理由＋附加好处｜⑥—⑧ 交还决定权}\par

\vspace{6pt}
{\color{accent}\bfseries 二、真题落地（2023 真题 · 艺术展还是机器人展）}\par
\vspace{2pt}
\noindent\fcolorbox{sandline}{sand}{\begin{minipage}{\dimexpr\linewidth-2\fboxsep-2\fboxrule\relax}\vspace{3pt}\small\linespread{1.03}\selectfont Dear David,\par
\vspace{3pt}
{\color{accent}\bfseries ①}You asked me for my advice on which of the two events is worth going to. {\color{accent}\bfseries ②}I would suggest the robot show.\par
\vspace{2pt}
{\color{accent}\bfseries ③}For one thing, it displays the latest home robots, which are close to our daily life. {\color{accent}\bfseries ④}For another, visitors can try operating a simple robot themselves. {\color{accent}\bfseries ⑤}Besides, the show is held indoors, which means the weather will not matter.\par
\vspace{2pt}
{\color{accent}\bfseries ⑥}Of course, the final choice is yours. {\color{accent}\bfseries ⑦}I have seen the poster, and it looks really attractive. {\color{accent}\bfseries ⑧}If you decide to go, I would be happy to go with you, and I am sure you will enjoy it.\par
\vspace{2pt}
\vspace{2pt}Yours sincerely,\par
Li Ming\par\vspace{3pt}\end{minipage}}\par

\vspace{6pt}
{\color{accent}\bfseries 三、换装表（只改这几处）}\par
\vspace{2pt}
{\footnotesize
\begin{tabularx}{\linewidth}{@{}>{\raggedright\arraybackslash}p{3.5cm}>{\raggedright\arraybackslash}p{5.6cm}X@{}}
\toprule
{\color{accent}\bfseries 槽位} & {\color{accent}\bfseries 本篇真题怎么填} & {\color{accent}\bfseries 可换词块} \\
\midrule
【建议对象】 & the robot show & the art exhibition / the online course \\
【理由一】【理由二】 & close to daily life / hands-on & free of charge / good for beginners \\
【陪同安排】 & go with you & help you book the ticket \\
\bottomrule
\end{tabularx}}\par

\vspace{5pt}
\noindent\begin{minipage}[t]{0.555\linewidth}
{\color{accent}\bfseries 四、阅卷人找什么}\par
{\footnotesize\begin{itemize}[leftmargin=1.2em,itemsep=1pt,topsep=1pt,parsep=0pt]\setlength{\leftmargin}{1.1em}
\item 第一段就给出明确选择，不写“both are good”这种和稀泥
\item 两条理由不重复：一条讲内容，一条讲体验或成本
\item 结尾把决定权交还对方，语气不命令
\end{itemize}}
\end{minipage}\hfill
\begin{minipage}[t]{0.415\linewidth}
{\color{hot}\bfseries 别踩这两个坑}\par
{\footnotesize\begin{itemize}[leftmargin=1.2em,itemsep=1pt,topsep=1pt,parsep=0pt]\setlength{\leftmargin}{1.1em}
\item 只给建议不给理由 —— 2023、2024 都明确要求 reason(s)
\item 两条理由其实是一条（都说 interesting），等于只有一条
\end{itemize}}
\end{minipage}\par

\vspace{6pt}
{\color{accent}\bfseries 五、想上 9 分，从这两句里挑一句加在结尾}\par
\vspace{2pt}
{\footnotesize
\begin{itemize}[leftmargin=1.2em,itemsep=1.5pt,topsep=1pt,parsep=0pt]
\item What makes it stand out is the hands-on experience it offers.\quad{\color{gray}（what 从句强调，替代 it is very good）}
\item Either choice would be fine, but this one fits your interest better.\quad{\color{gray}（让步＋比较，语气成熟，不武断）}
\end{itemize}}\par

\vspace{6pt}
{\color{accent}\bfseries 六、题干变体怎么改（同一篇，换两个词）}\par
\vspace{2pt}
{\footnotesize
\begin{itemize}[leftmargin=1.2em,itemsep=1.5pt,topsep=1pt,parsep=0pt]
\item 2019 辩论：建议对象换成 a specific topic，理由换成 relevance 与 data。
\item 2024 调查计划：建议对象换成 the plan itself，理由讲可行性与时间。
\end{itemize}}\par


\clearpage
\phantomsection\label{g05}
\noindent{\color{accent}\rule{\linewidth}{1.5pt}}\par\vspace{3pt}
\noindent{\fontsize{26}{28}\selectfont\bfseries\color{accent}05}\hspace{6pt}
{\LARGE\bfseries 感谢信}\hspace{7pt}{\large\color{gray} Thanks}\hfill
{\footnotesize\color{gray} 真题 2010 · 2016}\par
\vspace{1pt}
{\color{accent!35}\rule{\linewidth}{0.5pt}}\par\vspace{3pt}
{\footnotesize {\color{accent}\bfseries 题干信号}\quad 题干出现 thank / grateful / 感谢对方的接待、帮助或祝贺}\par
{\footnotesize {\color{accent}\bfseries 写作身份}\quad 受助者；对师长正式，对朋友可暖一些}\par
{\footnotesize {\color{accent}\bfseries 必须写到}\quad ①感谢的具体事项　②一处细节回忆　③这件事对你的影响　④回请或祝福}\par
{\footnotesize {\color{accent}\bfseries 一句话策略}\quad 不要只说 thank you——用“一个细节”证明感谢是真的}\par

\vspace{5pt}
{\color{accent}\bfseries 一、整篇模板（背这一篇）}\hfill{\footnotesize\color{gray} 正文 103 词 · 7 句}\par
\vspace{2pt}
\noindent\fcolorbox{rule}{soft}{\begin{minipage}{\dimexpr\linewidth-2\fboxsep-2\fboxrule\relax}\vspace{3pt}\small\linespread{1.03}\selectfont Dear {\color{hot}\bfseries 【称呼】},\par
\vspace{3pt}
{\color{accent}\bfseries ①}I have been back in {\color{hot}\bfseries 【地点】} for a week, but I still think of the days I spent in {\color{hot}\bfseries 【对方城市】}.{\color{accent}\bfseries ②}I am writing to thank you for {\color{hot}\bfseries 【感谢事项】}.\par
\vspace{2pt}
{\color{accent}\bfseries ③}Without your help, {\color{hot}\bfseries 【假设后果】}.{\color{accent}\bfseries ④}What impressed me most was {\color{hot}\bfseries 【具体细节】}.{\color{accent}\bfseries ⑤}I have told my family about it, and they all hope to meet you one day.\par
\vspace{2pt}
{\color{accent}\bfseries ⑥}By the way, you are welcome to visit {\color{hot}\bfseries 【我的城市】} at any time.{\color{accent}\bfseries ⑦}I would be glad to {\color{hot}\bfseries 【接待安排】} and introduce you to my family.\par
\vspace{2pt}
\vspace{2pt}Yours sincerely,\par
Li Ming\par\vspace{3pt}\end{minipage}}\par
\vspace{1pt}
{\footnotesize\color{gray} 段落分工｜①—② 背景＋感谢｜③—⑤ 影响＋细节｜⑥—⑦ 回请}\par

\vspace{6pt}
{\color{accent}\bfseries 二、真题落地（2010 真题 · 感谢接待并欢迎访华）}\par
\vspace{2pt}
\noindent\fcolorbox{sandline}{sand}{\begin{minipage}{\dimexpr\linewidth-2\fboxsep-2\fboxrule\relax}\vspace{3pt}\small\linespread{1.03}\selectfont Dear Mr. Smith,\par
\vspace{3pt}
{\color{accent}\bfseries ①}I have been back in Beijing for a week, but I still think of the days I spent in New York. {\color{accent}\bfseries ②}I am writing to thank you for your warm reception during my stay.\par
\vspace{2pt}
{\color{accent}\bfseries ③}Without your help, I could hardly have settled down so quickly. {\color{accent}\bfseries ④}What impressed me most was the dinner you prepared for me on my first night. {\color{accent}\bfseries ⑤}I have told my family about it, and they all hope to meet you one day.\par
\vspace{2pt}
{\color{accent}\bfseries ⑥}By the way, you are welcome to visit China at any time. {\color{accent}\bfseries ⑦}I would be glad to show you around Beijing and introduce you to my family.\par
\vspace{2pt}
\vspace{2pt}Yours sincerely,\par
Zhang Wei\par\vspace{3pt}\end{minipage}}\par

\vspace{6pt}
{\color{accent}\bfseries 三、换装表（只改这几处）}\par
\vspace{2pt}
{\footnotesize
\begin{tabularx}{\linewidth}{@{}>{\raggedright\arraybackslash}p{3.5cm}>{\raggedright\arraybackslash}p{5.6cm}X@{}}
\toprule
{\color{accent}\bfseries 槽位} & {\color{accent}\bfseries 本篇真题怎么填} & {\color{accent}\bfseries 可换词块} \\
\midrule
【感谢事项】 & your warm reception during my stay & your help with my application \\
【具体细节】 & the dinner on my first night & the books you lent me \\
【接待安排】 & show you around Beijing & take you to my hometown \\
\bottomrule
\end{tabularx}}\par

\vspace{5pt}
\noindent\begin{minipage}[t]{0.555\linewidth}
{\color{accent}\bfseries 四、阅卷人找什么}\par
{\footnotesize\begin{itemize}[leftmargin=1.2em,itemsep=1pt,topsep=1pt,parsep=0pt]\setlength{\leftmargin}{1.1em}
\item 感谢落到具体事上（接待、修改文书、借书），不是空洞的 your help
\item 有一处细节回忆，这是 8 分以上的分水岭
\item 结尾有回请或祝福，完成“礼尚往来”
\end{itemize}}
\end{minipage}\hfill
\begin{minipage}[t]{0.415\linewidth}
{\color{hot}\bfseries 别踩这两个坑}\par
{\footnotesize\begin{itemize}[leftmargin=1.2em,itemsep=1pt,topsep=1pt,parsep=0pt]\setlength{\leftmargin}{1.1em}
\item 全篇 thank you so much，没有一件具体的事
\item 感谢信用了过去时又突然跳到 I will，时态来回跳
\end{itemize}}
\end{minipage}\par

\vspace{6pt}
{\color{accent}\bfseries 五、想上 9 分，从这两句里挑一句加在结尾}\par
\vspace{2pt}
{\footnotesize
\begin{itemize}[leftmargin=1.2em,itemsep=1.5pt,topsep=1pt,parsep=0pt]
\item Your kindness made a stranger feel at home.\quad{\color{gray}（make sb. feel at home，比 warm reception 更有画面）}
\item I hope I can return the favour when you come to China.\quad{\color{gray}（return the favour 是“还人情”的地道说法）}
\end{itemize}}\par

\vspace{6pt}
{\color{accent}\bfseries 六、题干变体怎么改（同一篇，换两个词）}\par
\vspace{2pt}
{\footnotesize
\begin{itemize}[leftmargin=1.2em,itemsep=1.5pt,topsep=1pt,parsep=0pt]
\item 2016 回复祝贺：开头改 Thank you for your congratulations，再接建议。
\item 感谢对象是机构时，称呼用 Dear Sir or Madam，落款仍用题目姓名。
\end{itemize}}\par


\clearpage
\phantomsection\label{g06}
\noindent{\color{accent}\rule{\linewidth}{1.5pt}}\par\vspace{3pt}
\noindent{\fontsize{26}{28}\selectfont\bfseries\color{accent}06}\hspace{6pt}
{\LARGE\bfseries 祝贺信}\hspace{7pt}{\large\color{gray} Congratulations}\hfill
{\footnotesize\color{gray} 真题 2011 · 2016（回复祝贺）}\par
\vspace{1pt}
{\color{accent!35}\rule{\linewidth}{0.5pt}}\par\vspace{3pt}
{\footnotesize {\color{accent}\bfseries 题干信号}\quad 题干出现 congratulate / 祝贺对方录取、获奖、晋升}\par
{\footnotesize {\color{accent}\bfseries 写作身份}\quad 亲戚、朋友或同学；语气热情但别夸张}\par
{\footnotesize {\color{accent}\bfseries 必须写到}\quad ①祝贺的具体事　②为何值得祝贺　③两条建议或祝福　④保持联系}\par
{\footnotesize {\color{accent}\bfseries 一句话策略}\quad 祝贺一句到位，第二段立刻转到“接下来怎么做”，避免整篇客套}\par

\vspace{5pt}
{\color{accent}\bfseries 一、整篇模板（背这一篇）}\hfill{\footnotesize\color{gray} 正文 96 词 · 9 句}\par
\vspace{2pt}
\noindent\fcolorbox{rule}{soft}{\begin{minipage}{\dimexpr\linewidth-2\fboxsep-2\fboxrule\relax}\vspace{3pt}\small\linespread{1.03}\selectfont Dear {\color{hot}\bfseries 【称呼】},\par
\vspace{3pt}
{\color{accent}\bfseries ①}I was delighted to learn that {\color{hot}\bfseries 【祝贺事项】}.{\color{accent}\bfseries ②}Congratulations on {\color{hot}\bfseries 【具体成就】}!\par
\vspace{2pt}
{\color{accent}\bfseries ③}As you are about to {\color{hot}\bfseries 【阶段】}, I would like to offer two suggestions.{\color{accent}\bfseries ④}First, {\color{hot}\bfseries 【建议一】}.{\color{accent}\bfseries ⑤}Second, {\color{hot}\bfseries 【建议二】}.{\color{accent}\bfseries ⑥}These are only my personal views, and you know your own situation best.\par
\vspace{2pt}
{\color{accent}\bfseries ⑦}I am sure you will do well.{\color{accent}\bfseries ⑧}I wish you a happy and fruitful four years.{\color{accent}\bfseries ⑨}Please keep in touch.\par
\vspace{2pt}
\vspace{2pt}Yours sincerely,\par
Li Ming\par\vspace{3pt}\end{minipage}}\par
\vspace{1pt}
{\footnotesize\color{gray} 段落分工｜①—② 祝贺｜③—⑥ 两条建议｜⑦—⑨ 祝福＋联系}\par

\vspace{6pt}
{\color{accent}\bfseries 二、真题落地（2011 真题 · 祝贺录取并给大学适应建议）}\par
\vspace{2pt}
\noindent\fcolorbox{sandline}{sand}{\begin{minipage}{\dimexpr\linewidth-2\fboxsep-2\fboxrule\relax}\vspace{3pt}\small\linespread{1.03}\selectfont Dear Li Ming,\par
\vspace{3pt}
{\color{accent}\bfseries ①}I was delighted to learn that you have been admitted to Beijing University. {\color{accent}\bfseries ②}Congratulations on this great achievement!\par
\vspace{2pt}
{\color{accent}\bfseries ③}As you are about to start university life, I would like to offer two suggestions. {\color{accent}\bfseries ④}First, join one or two clubs so that you can meet people outside your major. {\color{accent}\bfseries ⑤}Second, plan your time carefully, for no one will remind you of deadlines any more. {\color{accent}\bfseries ⑥}These are only my personal views, and you know your own situation best.\par
\vspace{2pt}
{\color{accent}\bfseries ⑦}I am sure you will do well. {\color{accent}\bfseries ⑧}I wish you a happy and fruitful four years. {\color{accent}\bfseries ⑨}Please keep in touch.\par
\vspace{2pt}
\vspace{2pt}Yours sincerely,\par
Zhang Wei\par\vspace{3pt}\end{minipage}}\par

\vspace{6pt}
{\color{accent}\bfseries 三、换装表（只改这几处）}\par
\vspace{2pt}
{\footnotesize
\begin{tabularx}{\linewidth}{@{}>{\raggedright\arraybackslash}p{3.5cm}>{\raggedright\arraybackslash}p{5.6cm}X@{}}
\toprule
{\color{accent}\bfseries 槽位} & {\color{accent}\bfseries 本篇真题怎么填} & {\color{accent}\bfseries 可换词块} \\
\midrule
【祝贺事项】 & you have been admitted to Beijing University & you won the translation contest \\
【阶段】 & start university life & begin your new job \\
【建议一】【建议二】 & join clubs / plan your time & ask questions in class / save part of your salary \\
\bottomrule
\end{tabularx}}\par

\vspace{5pt}
\noindent\begin{minipage}[t]{0.555\linewidth}
{\color{accent}\bfseries 四、阅卷人找什么}\par
{\footnotesize\begin{itemize}[leftmargin=1.2em,itemsep=1pt,topsep=1pt,parsep=0pt]\setlength{\leftmargin}{1.1em}
\item 祝贺的事写得具体（哪所学校、什么奖），不写 good news
\item 建议两条且不重复，每条带一点理由
\item 结尾有祝福，不写“祝你成功”式的中文直译
\end{itemize}}
\end{minipage}\hfill
\begin{minipage}[t]{0.415\linewidth}
{\color{hot}\bfseries 别踩这两个坑}\par
{\footnotesize\begin{itemize}[leftmargin=1.2em,itemsep=1pt,topsep=1pt,parsep=0pt]\setlength{\leftmargin}{1.1em}
\item Congratulations for → 必须用 on
\item 整篇只有祝贺，没有第二个动作（2011 还要求 suggestions）
\end{itemize}}
\end{minipage}\par

\vspace{6pt}
{\color{accent}\bfseries 五、想上 9 分，从这两句里挑一句加在结尾}\par
\vspace{2pt}
{\footnotesize
\begin{itemize}[leftmargin=1.2em,itemsep=1.5pt,topsep=1pt,parsep=0pt]
\item This is only the beginning of a much longer journey.\quad{\color{gray}（比喻收尾，比 Good luck 高一档）}
\item Do let me know how you settle in.\quad{\color{gray}（do 强调＋settle in，口语得体）}
\end{itemize}}\par

\vspace{6pt}
{\color{accent}\bfseries 六、题干变体怎么改（同一篇，换两个词）}\par
\vspace{2pt}
{\footnotesize
\begin{itemize}[leftmargin=1.2em,itemsep=1.5pt,topsep=1pt,parsep=0pt]
\item 2016 回复祝贺：第一段感谢对方的祝贺，第二段再给翻译建议。
\item 祝贺教授晋升时，语气收一点，不用感叹号连发。
\end{itemize}}\par


\clearpage
\phantomsection\label{g07}
\noindent{\color{accent}\rule{\linewidth}{1.5pt}}\par\vspace{3pt}
\noindent{\fontsize{26}{28}\selectfont\bfseries\color{accent}07}\hspace{6pt}
{\LARGE\bfseries 道歉信}\hspace{7pt}{\large\color{gray} Apology}\hfill
{\footnotesize\color{gray} 真题 2018}\par
\vspace{1pt}
{\color{accent!35}\rule{\linewidth}{0.5pt}}\par\vspace{3pt}
{\footnotesize {\color{accent}\bfseries 题干信号}\quad 题干出现 apologize / 取消、迟到、弄丢、不能赴约，并要求补救}\par
{\footnotesize {\color{accent}\bfseries 写作身份}\quad 失约方；对师长正式，对朋友可少一点公文腔}\par
{\footnotesize {\color{accent}\bfseries 必须写到}\quad ①为何道歉　②一句原因　③补救方案　④另约时间或替代办法}\par
{\footnotesize {\color{accent}\bfseries 一句话策略}\quad 道歉—原因—补救—另约，四句四件事，原因绝不超过一句}\par

\vspace{5pt}
{\color{accent}\bfseries 一、整篇模板（背这一篇）}\hfill{\footnotesize\color{gray} 正文 97 词 · 9 句}\par
\vspace{2pt}
\noindent\fcolorbox{rule}{soft}{\begin{minipage}{\dimexpr\linewidth-2\fboxsep-2\fboxrule\relax}\vspace{3pt}\small\linespread{1.03}\selectfont Dear {\color{hot}\bfseries 【称呼】},\par
\vspace{3pt}
{\color{accent}\bfseries ①}I am writing to apologize for {\color{hot}\bfseries 【道歉事项】}.{\color{accent}\bfseries ②}I am afraid I have to {\color{hot}\bfseries 【取消/推迟】} at such short notice.\par
\vspace{2pt}
{\color{accent}\bfseries ③}The reason is that {\color{hot}\bfseries 【原因】}.{\color{accent}\bfseries ④}I know this will cause you trouble, and I am truly sorry.{\color{accent}\bfseries ⑤}To make up for it, I would like to {\color{hot}\bfseries 【补救方案】}.{\color{accent}\bfseries ⑥}Would {\color{hot}\bfseries 【备选时间】} suit you?{\color{accent}\bfseries ⑦}I can also {\color{hot}\bfseries 【替代安排】} if that is more convenient.\par
\vspace{2pt}
{\color{accent}\bfseries ⑧}Please accept my apology.{\color{accent}\bfseries ⑨}I hope we can meet soon.\par
\vspace{2pt}
\vspace{2pt}Yours sincerely,\par
Li Ming\par\vspace{3pt}\end{minipage}}\par
\vspace{1pt}
{\footnotesize\color{gray} 段落分工｜①—② 道歉＋取消｜③—⑦ 原因＋补救＋另约｜⑧—⑨ 收尾}\par

\vspace{6pt}
{\color{accent}\bfseries 二、真题落地（2018 真题 · 取消拜访并另约）}\par
\vspace{2pt}
\noindent\fcolorbox{sandline}{sand}{\begin{minipage}{\dimexpr\linewidth-2\fboxsep-2\fboxrule\relax}\vspace{3pt}\small\linespread{1.03}\selectfont Dear Professor Smith,\par
\vspace{3pt}
{\color{accent}\bfseries ①}I am writing to apologize for cancelling my visit to you next Monday. {\color{accent}\bfseries ②}I am afraid I have to put it off at such short notice.\par
\vspace{2pt}
{\color{accent}\bfseries ③}The reason is that my mother has suddenly fallen ill and I must take her to hospital. {\color{accent}\bfseries ④}I know this will cause you trouble, and I am truly sorry. {\color{accent}\bfseries ⑤}To make up for it, I would like to visit you the following Friday. {\color{accent}\bfseries ⑥}Would that suit you? {\color{accent}\bfseries ⑦}I can also send you my questions by email if that is more convenient.\par
\vspace{2pt}
{\color{accent}\bfseries ⑧}Please accept my apology. {\color{accent}\bfseries ⑨}I hope we can meet soon.\par
\vspace{2pt}
\vspace{2pt}Yours sincerely,\par
Li Ming\par\vspace{3pt}\end{minipage}}\par

\vspace{6pt}
{\color{accent}\bfseries 三、换装表（只改这几处）}\par
\vspace{2pt}
{\footnotesize
\begin{tabularx}{\linewidth}{@{}>{\raggedright\arraybackslash}p{3.5cm}>{\raggedright\arraybackslash}p{5.6cm}X@{}}
\toprule
{\color{accent}\bfseries 槽位} & {\color{accent}\bfseries 本篇真题怎么填} & {\color{accent}\bfseries 可换词块} \\
\midrule
【道歉事项】 & cancelling my visit to you next Monday & missing your lecture yesterday \\
【原因】 & my mother has suddenly fallen ill & I was sent on an urgent business trip \\
【补救方案】 & visit you the following Friday & call you this weekend \\
\bottomrule
\end{tabularx}}\par

\vspace{5pt}
\noindent\begin{minipage}[t]{0.555\linewidth}
{\color{accent}\bfseries 四、阅卷人找什么}\par
{\footnotesize\begin{itemize}[leftmargin=1.2em,itemsep=1pt,topsep=1pt,parsep=0pt]\setlength{\leftmargin}{1.1em}
\item 道歉信必须同时有“原因”和“补救”，只有 sorry 最多 5 分
\item 原因一句收住，不写成记叙文
\item 另约时间具体（the following Friday），不用 someday / later
\end{itemize}}
\end{minipage}\hfill
\begin{minipage}[t]{0.415\linewidth}
{\color{hot}\bfseries 别踩这两个坑}\par
{\footnotesize\begin{itemize}[leftmargin=1.2em,itemsep=1pt,topsep=1pt,parsep=0pt]\setlength{\leftmargin}{1.1em}
\item 把道歉写成解释，通篇 because，最后没有补救
\item 对教授用太随意的 Sorry about that!，语域不合
\end{itemize}}
\end{minipage}\par

\vspace{6pt}
{\color{accent}\bfseries 五、想上 9 分，从这两句里挑一句加在结尾}\par
\vspace{2pt}
{\footnotesize
\begin{itemize}[leftmargin=1.2em,itemsep=1.5pt,topsep=1pt,parsep=0pt]
\item I fully understand how inconvenient this must be for you.\quad{\color{gray}（先共情再解释，比单说 sorry 成熟）}
\item I will make sure it does not happen again.\quad{\color{gray}（承诺具体，补救不止一句空话）}
\end{itemize}}\par

\vspace{6pt}
{\color{accent}\bfseries 六、题干变体怎么改（同一篇，换两个词）}\par
\vspace{2pt}
{\footnotesize
\begin{itemize}[leftmargin=1.2em,itemsep=1.5pt,topsep=1pt,parsep=0pt]
\item 2018 取消拜访；若只是迟到，把取消改成迟到，补救改成改期。
\item 对朋友可删掉 Please accept my apology，换成 I owe you one。
\end{itemize}}\par


\clearpage
\phantomsection\label{g08}
\noindent{\color{accent}\rule{\linewidth}{1.5pt}}\par\vspace{3pt}
\noindent{\fontsize{26}{28}\selectfont\bfseries\color{accent}08}\hspace{6pt}
{\LARGE\bfseries 投诉信}\hspace{7pt}{\large\color{gray} Complaint}\hfill
{\footnotesize\color{gray} 真题 2012}\par
\vspace{1pt}
{\color{accent!35}\rule{\linewidth}{0.5pt}}\par\vspace{3pt}
{\footnotesize {\color{accent}\bfseries 题干信号}\quad 题干出现 complain / 商品或服务有问题，要求解决或赔偿}\par
{\footnotesize {\color{accent}\bfseries 写作身份}\quad 消费者；对商家用 Sir or Madam，对具体人用其姓氏}\par
{\footnotesize {\color{accent}\bfseries 必须写到}\quad ①购买渠道与时间　②具体故障　③已尝试的解决办法　④明确要求与期限}\par
{\footnotesize {\color{accent}\bfseries 一句话策略}\quad 事实—故障—影响—要求，全程冷静，不写情绪词}\par

\vspace{5pt}
{\color{accent}\bfseries 一、整篇模板（背这一篇）}\hfill{\footnotesize\color{gray} 正文 95 词 · 7 句}\par
\vspace{2pt}
\noindent\fcolorbox{rule}{soft}{\begin{minipage}{\dimexpr\linewidth-2\fboxsep-2\fboxrule\relax}\vspace{3pt}\small\linespread{1.03}\selectfont Dear {\color{hot}\bfseries 【称呼】},\par
\vspace{3pt}
{\color{accent}\bfseries ①}I am writing to complain about {\color{hot}\bfseries 【商品】}, which I bought from {\color{hot}\bfseries 【购买渠道】} {\color{hot}\bfseries 【购买日期】}.\par
\vspace{2pt}
{\color{accent}\bfseries ②}The problem is that {\color{hot}\bfseries 【问题】}.{\color{accent}\bfseries ③}I have tried {\color{hot}\bfseries 【已采取措施】}, but the fault remains.{\color{accent}\bfseries ④}This has made it impossible for me to {\color{hot}\bfseries 【受影响的事】}.\par
\vspace{2pt}
{\color{accent}\bfseries ⑤}I would like you to {\color{hot}\bfseries 【要求一】}, and I hope you can {\color{hot}\bfseries 【要求二】} within {\color{hot}\bfseries 【期限】}.{\color{accent}\bfseries ⑥}If not, I shall have to {\color{hot}\bfseries 【进一步措施】}.{\color{accent}\bfseries ⑦}I look forward to your reply.\par
\vspace{2pt}
\vspace{2pt}Yours sincerely,\par
Li Ming\par\vspace{3pt}\end{minipage}}\par
\vspace{1pt}
{\footnotesize\color{gray} 段落分工｜① 购买信息｜②—④ 故障＋影响｜⑤—⑦ 要求＋期限}\par

\vspace{6pt}
{\color{accent}\bfseries 二、真题落地（2012 真题 · 电子词典故障）}\par
\vspace{2pt}
\noindent\fcolorbox{sandline}{sand}{\begin{minipage}{\dimexpr\linewidth-2\fboxsep-2\fboxrule\relax}\vspace{3pt}\small\linespread{1.03}\selectfont Dear Sir or Madam,\par
\vspace{3pt}
{\color{accent}\bfseries ①}I am writing to complain about the electronic dictionary, which I bought from your online store last Friday.\par
\vspace{2pt}
{\color{accent}\bfseries ②}The problem is that the screen goes black whenever I turn it on. {\color{accent}\bfseries ③}I have tried changing the batteries and resetting it, but the fault remains. {\color{accent}\bfseries ④}This has made it impossible for me to prepare for my exam.\par
\vspace{2pt}
{\color{accent}\bfseries ⑤}I would like you to replace it with a new one, and I hope you can arrange the delivery within a week. {\color{accent}\bfseries ⑥}If not, I shall have to ask for a full refund. {\color{accent}\bfseries ⑦}I look forward to your reply.\par
\vspace{2pt}
\vspace{2pt}Yours sincerely,\par
Zhang Wei\par\vspace{3pt}\end{minipage}}\par

\vspace{6pt}
{\color{accent}\bfseries 三、换装表（只改这几处）}\par
\vspace{2pt}
{\footnotesize
\begin{tabularx}{\linewidth}{@{}>{\raggedright\arraybackslash}p{3.5cm}>{\raggedright\arraybackslash}p{5.6cm}X@{}}
\toprule
{\color{accent}\bfseries 槽位} & {\color{accent}\bfseries 本篇真题怎么填} & {\color{accent}\bfseries 可换词块} \\
\midrule
【商品】 & the electronic dictionary & the laptop / the online course \\
【问题】 & the screen goes black & the battery lasts only an hour \\
【要求一】【要求二】 & replace it / deliver within a week & repair it free / refund the shipping fee \\
\bottomrule
\end{tabularx}}\par

\vspace{5pt}
\noindent\begin{minipage}[t]{0.555\linewidth}
{\color{accent}\bfseries 四、阅卷人找什么}\par
{\footnotesize\begin{itemize}[leftmargin=1.2em,itemsep=1pt,topsep=1pt,parsep=0pt]\setlength{\leftmargin}{1.1em}
\item 购买渠道＋日期写清楚，这是投诉信的身份证明
\item 故障描述具体到“什么时候、什么现象”，不写 it is broken
\item 要求明确（换、修、退），并给期限
\end{itemize}}
\end{minipage}\hfill
\begin{minipage}[t]{0.415\linewidth}
{\color{hot}\bfseries 别踩这两个坑}\par
{\footnotesize\begin{itemize}[leftmargin=1.2em,itemsep=1pt,topsep=1pt,parsep=0pt]\setlength{\leftmargin}{1.1em}
\item 通篇 angry / terrible / cheat，情绪代替事实
\item 只投诉不要求，任务第二个动作缺失
\end{itemize}}
\end{minipage}\par

\vspace{6pt}
{\color{accent}\bfseries 五、想上 9 分，从这两句里挑一句加在结尾}\par
\vspace{2pt}
{\footnotesize
\begin{itemize}[leftmargin=1.2em,itemsep=1.5pt,topsep=1pt,parsep=0pt]
\item I have kept the receipt and can send you a photo if needed.\quad{\color{gray}（证据意识，投诉可信度立刻上来）}
\item I would appreciate a reply at your earliest convenience.\quad{\color{gray}（礼貌催复，不写 I want a reply now）}
\end{itemize}}\par

\vspace{6pt}
{\color{accent}\bfseries 六、题干变体怎么改（同一篇，换两个词）}\par
\vspace{2pt}
{\footnotesize
\begin{itemize}[leftmargin=1.2em,itemsep=1.5pt,topsep=1pt,parsep=0pt]
\item 投诉服务（课程、快递）时，商品换成 the course / the delivery。
\item 要求部分按“换—修—退”顺序写，期限给具体天数。
\end{itemize}}\par


\clearpage
\phantomsection\label{g09}
\noindent{\color{accent}\rule{\linewidth}{1.5pt}}\par\vspace{3pt}
\noindent{\fontsize{26}{28}\selectfont\bfseries\color{accent}09}\hspace{6pt}
{\LARGE\bfseries 咨询信}\hspace{7pt}{\large\color{gray} Inquiry}\hfill
{\footnotesize\color{gray} 真题 2014}\par
\vspace{1pt}
{\color{accent!35}\rule{\linewidth}{0.5pt}}\par\vspace{3pt}
{\footnotesize {\color{accent}\bfseries 题干信号}\quad 题干出现 ask about / 询问信息，或“告知自己情况并征求意见”}\par
{\footnotesize {\color{accent}\bfseries 写作身份}\quad 即将前往当地的学生或准参与者；对当地人客气}\par
{\footnotesize {\color{accent}\bfseries 必须写到}\quad ①自我介绍或背景　②两到三个具体问题　③为什么问　④致谢}\par
{\footnotesize {\color{accent}\bfseries 一句话策略}\quad 问题要具体到“能不能、在哪里、多少钱”，别问“能介绍一下吗”}\par

\vspace{5pt}
{\color{accent}\bfseries 一、整篇模板（背这一篇）}\hfill{\footnotesize\color{gray} 正文 99 词 · 8 句}\par
\vspace{2pt}
\noindent\fcolorbox{rule}{soft}{\begin{minipage}{\dimexpr\linewidth-2\fboxsep-2\fboxrule\relax}\vspace{3pt}\small\linespread{1.03}\selectfont Dear {\color{hot}\bfseries 【称呼】},\par
\vspace{3pt}
{\color{accent}\bfseries ①}I am writing to ask for your advice about {\color{hot}\bfseries 【咨询主题】}.{\color{accent}\bfseries ②}I will move to {\color{hot}\bfseries 【地点】} next month, and I know little about {\color{hot}\bfseries 【当地情况】}.\par
\vspace{2pt}
{\color{accent}\bfseries ③}To give you some idea of my habits, I {\color{hot}\bfseries 【习惯一】}, and I usually {\color{hot}\bfseries 【习惯二】}.{\color{accent}\bfseries ④}I wonder whether {\color{hot}\bfseries 【问题一】}.{\color{accent}\bfseries ⑤}Besides, could you tell me {\color{hot}\bfseries 【问题二】}?\par
\vspace{2pt}
{\color{accent}\bfseries ⑥}I am a little worried about the first few weeks, as everything will be new to me.{\color{accent}\bfseries ⑦}Your suggestions would be of great help.{\color{accent}\bfseries ⑧}Thank you in advance.\par
\vspace{2pt}
\vspace{2pt}Yours sincerely,\par
Li Ming\par\vspace{3pt}\end{minipage}}\par
\vspace{1pt}
{\footnotesize\color{gray} 段落分工｜①—② 背景｜③—⑤ 习惯＋两个问题｜⑥—⑧ 致谢}\par

\vspace{6pt}
{\color{accent}\bfseries 二、真题落地（2014 真题 · 告知生活习惯并咨询合租）}\par
\vspace{2pt}
\noindent\fcolorbox{sandline}{sand}{\begin{minipage}{\dimexpr\linewidth-2\fboxsep-2\fboxrule\relax}\vspace{3pt}\small\linespread{1.03}\selectfont Dear John,\par
\vspace{3pt}
{\color{accent}\bfseries ①}I am writing to ask for your advice about daily life there. {\color{accent}\bfseries ②}I will move to London next month, and I know little about local customs.\par
\vspace{2pt}
{\color{accent}\bfseries ③}To give you some idea of my habits, I get up early and go to bed before eleven, and I usually cook my own meals. {\color{accent}\bfseries ④}I wonder whether my schedule will disturb you. {\color{accent}\bfseries ⑤}Besides, could you tell me where I can buy fresh vegetables nearby?\par
\vspace{2pt}
{\color{accent}\bfseries ⑥}I am a little worried about the first few weeks, as everything will be new to me. {\color{accent}\bfseries ⑦}Your suggestions would be of great help. {\color{accent}\bfseries ⑧}Thank you in advance.\par
\vspace{2pt}
\vspace{2pt}Yours sincerely,\par
Li Ming\par\vspace{3pt}\end{minipage}}\par

\vspace{6pt}
{\color{accent}\bfseries 三、换装表（只改这几处）}\par
\vspace{2pt}
{\footnotesize
\begin{tabularx}{\linewidth}{@{}>{\raggedright\arraybackslash}p{3.5cm}>{\raggedright\arraybackslash}p{5.6cm}X@{}}
\toprule
{\color{accent}\bfseries 槽位} & {\color{accent}\bfseries 本篇真题怎么填} & {\color{accent}\bfseries 可换词块} \\
\midrule
【咨询主题】 & daily life there & the language course / the apartment \\
【问题一】【问题二】 & will my schedule disturb you / where to buy food & is the deposit refundable / when does the course start \\
【当地情况】 & local customs & the transport system \\
\bottomrule
\end{tabularx}}\par

\vspace{5pt}
\noindent\begin{minipage}[t]{0.555\linewidth}
{\color{accent}\bfseries 四、阅卷人找什么}\par
{\footnotesize\begin{itemize}[leftmargin=1.2em,itemsep=1pt,topsep=1pt,parsep=0pt]\setlength{\leftmargin}{1.1em}
\item 问题两到三个，且每个都能用一句话回答（不是开放式的“介绍一下”）
\item 先给背景再提问，让对方知道你的情况
\item 结尾致谢，不写“等待您的答复”式的中文直译
\end{itemize}}
\end{minipage}\hfill
\begin{minipage}[t]{0.415\linewidth}
{\color{hot}\bfseries 别踩这两个坑}\par
{\footnotesize\begin{itemize}[leftmargin=1.2em,itemsep=1pt,topsep=1pt,parsep=0pt]\setlength{\leftmargin}{1.1em}
\item 只写 I want to know more，没有一个具体问题
\item 把咨询信写成自我介绍，通篇讲自己，问题只有一句
\end{itemize}}
\end{minipage}\par

\vspace{6pt}
{\color{accent}\bfseries 五、想上 9 分，从这两句里挑一句加在结尾}\par
\vspace{2pt}
{\footnotesize
\begin{itemize}[leftmargin=1.2em,itemsep=1.5pt,topsep=1pt,parsep=0pt]
\item Any advice you could offer would be much appreciated.\quad{\color{gray}（被动礼貌请求，比 please tell me 委婉）}
\item I would also be grateful for a rough idea of the cost.\quad{\color{gray}（把问题问得更具体，对方好回答）}
\end{itemize}}\par

\vspace{6pt}
{\color{accent}\bfseries 六、题干变体怎么改（同一篇，换两个词）}\par
\vspace{2pt}
{\footnotesize
\begin{itemize}[leftmargin=1.2em,itemsep=1.5pt,topsep=1pt,parsep=0pt]
\item 2024 征求意见：问题改成 Do you think the plan is workable?
\item 咨询住宿时，两个问题分别问价格与距离，不要都问好不好。
\end{itemize}}\par


\clearpage
\phantomsection\label{g10}
\noindent{\color{accent}\rule{\linewidth}{1.5pt}}\par\vspace{3pt}
\noindent{\fontsize{26}{28}\selectfont\bfseries\color{accent}10}\hspace{6pt}
{\LARGE\bfseries 通知}\hspace{7pt}{\large\color{gray} Notice}\hfill
{\footnotesize\color{gray} 真题 2013 · 2015}\par
\vspace{1pt}
{\color{accent!35}\rule{\linewidth}{0.5pt}}\par\vspace{3pt}
{\footnotesize {\color{accent}\bfseries 题干信号}\quad 题干出现 Notice / 通知同学参加，或“介绍活动并招募志愿者”}\par
{\footnotesize {\color{accent}\bfseries 写作身份}\quad 学生会、社团或班委；以集体名义发布，不写个人称呼}\par
{\footnotesize {\color{accent}\bfseries 必须写到}\quad ①标题　②活动名称、时间、地点　③活动内容　④报名方式或招募要求}\par
{\footnotesize {\color{accent}\bfseries 一句话策略}\quad 通知没有称呼；标题居中；落款单位＋日期在右下角}\par

\vspace{5pt}
{\color{accent}\bfseries 一、整篇模板（背这一篇）}\hfill{\footnotesize\color{gray} 正文 96 词 · 8 句}\par
\vspace{2pt}
\noindent\fcolorbox{rule}{soft}{\begin{minipage}{\dimexpr\linewidth-2\fboxsep-2\fboxrule\relax}\vspace{3pt}\small\linespread{1.03}\selectfont {\centering\bfseries\large Notice\par}
\vspace{2pt}
{\color{accent}\bfseries ①}{\color{hot}\bfseries 【活动名称】} will be held {\color{hot}\bfseries 【时间】} at {\color{hot}\bfseries 【地点】}.{\color{accent}\bfseries ②}It is organized by {\color{hot}\bfseries 【主办方】} for {\color{hot}\bfseries 【参加对象】}.\par
\vspace{2pt}
{\color{accent}\bfseries ③}The program includes {\color{hot}\bfseries 【活动一】}, {\color{hot}\bfseries 【活动二】} and {\color{hot}\bfseries 【活动三】}.{\color{accent}\bfseries ④}{\color{hot}\bfseries 【附加安排】}.{\color{accent}\bfseries ⑤}Those who wish to take part should {\color{hot}\bfseries 【报名方式】} before {\color{hot}\bfseries 【截止日期】}.\par
\vspace{2pt}
{\color{accent}\bfseries ⑥}We are also looking for {\color{hot}\bfseries 【招募人数】} volunteers to {\color{hot}\bfseries 【志愿工作】}.{\color{accent}\bfseries ⑦}Anyone interested may contact {\color{hot}\bfseries 【联系人】} at {\color{hot}\bfseries 【联系方式】}.{\color{accent}\bfseries ⑧}Please apply early, as places are limited.\par
\vspace{2pt}
\begin{flushright}{\color{hot}\bfseries 【主办方】}\\{\color{hot}\bfseries 【日期】}\end{flushright}\vspace{3pt}\end{minipage}}\par
\vspace{1pt}
{\footnotesize\color{gray} 段落分工｜①—② 活动概况｜③—⑤ 内容＋报名｜⑥—⑧ 招募＋截止}\par

\vspace{6pt}
{\color{accent}\bfseries 二、真题落地（2015 真题 · 夏令营介绍＋招募志愿者）}\par
\vspace{2pt}
\noindent\fcolorbox{sandline}{sand}{\begin{minipage}{\dimexpr\linewidth-2\fboxsep-2\fboxrule\relax}\vspace{3pt}\small\linespread{1.03}\selectfont {\centering\bfseries\large Notice\par}
\vspace{2pt}
{\color{accent}\bfseries ①}The Summer Camp for high school students will be held from July 10 to July 20 at our campus. {\color{accent}\bfseries ②}It is organized by the Student Union for local high school students.\par
\vspace{2pt}
{\color{accent}\bfseries ③}The program includes campus tours, Chinese culture classes and a city walk. {\color{accent}\bfseries ④}Campers will stay in the student dormitory. {\color{accent}\bfseries ⑤}Those who wish to take part should sign up at the Student Union office before July 1.\par
\vspace{2pt}
{\color{accent}\bfseries ⑥}We are also looking for twenty volunteers to help with reception and translation. {\color{accent}\bfseries ⑦}Anyone interested may contact Li Ming at Room 203. {\color{accent}\bfseries ⑧}Please apply early, as places are limited.\par
\vspace{2pt}
\begin{flushright}The Student Union\\June 20, 2026\end{flushright}\vspace{3pt}\end{minipage}}\par

\vspace{6pt}
{\color{accent}\bfseries 三、换装表（只改这几处）}\par
\vspace{2pt}
{\footnotesize
\begin{tabularx}{\linewidth}{@{}>{\raggedright\arraybackslash}p{3.5cm}>{\raggedright\arraybackslash}p{5.6cm}X@{}}
\toprule
{\color{accent}\bfseries 槽位} & {\color{accent}\bfseries 本篇真题怎么填} & {\color{accent}\bfseries 可换词块} \\
\midrule
【活动名称】 & The Summer Camp & A Charity Sale / A Campus Concert \\
【报名方式】 & sign up at the Student Union office & send an email to notice@univ.edu \\
【志愿工作】 & help with reception and translation & set up the stage / guide the visitors \\
【附加安排】 & Campers will stay in the student dormitory & Lunch will be provided / Donations are welcome \\
\bottomrule
\end{tabularx}}\par

\vspace{5pt}
\noindent\begin{minipage}[t]{0.555\linewidth}
{\color{accent}\bfseries 四、阅卷人找什么}\par
{\footnotesize\begin{itemize}[leftmargin=1.2em,itemsep=1pt,topsep=1pt,parsep=0pt]\setlength{\leftmargin}{1.1em}
\item 标题 Notice 单独居中，落款单位＋日期在右下角
\item 活动信息（名称、时间、地点、内容）齐全
\item 招募类通知要写清人数、工作内容、报名方式与截止日
\end{itemize}}
\end{minipage}\hfill
\begin{minipage}[t]{0.415\linewidth}
{\color{hot}\bfseries 别踩这两个坑}\par
{\footnotesize\begin{itemize}[leftmargin=1.2em,itemsep=1pt,topsep=1pt,parsep=0pt]\setlength{\leftmargin}{1.1em}
\item 通知加了 Dear classmates, —— 格式直接降档
\item 只介绍活动不招募，或只招募不介绍（2015 两个动作）
\end{itemize}}
\end{minipage}\par

\vspace{6pt}
{\color{accent}\bfseries 五、想上 9 分，从这两句里挑一句加在结尾}\par
\vspace{2pt}
{\footnotesize
\begin{itemize}[leftmargin=1.2em,itemsep=1.5pt,topsep=1pt,parsep=0pt]
\item Places will be allocated on a first-come, first-served basis.\quad{\color{gray}（先到先得，通知里一句话说清规则）}
\item For further information, please visit our website.\quad{\color{gray}（通知常用收尾，替掉 Thank you for listening）}
\end{itemize}}\par

\vspace{6pt}
{\color{accent}\bfseries 六、题干变体怎么改（同一篇，换两个词）}\par
\vspace{2pt}
{\footnotesize
\begin{itemize}[leftmargin=1.2em,itemsep=1.5pt,topsep=1pt,parsep=0pt]
\item 2013 慈善义卖：标题仍 Notice，第三段改成鼓励参加，不写招募。
\item 线上活动把地点改 on the school website，报名方式给链接。
\end{itemize}}\par


\clearpage
\phantomsection\label{g11}
\noindent{\color{accent}\rule{\linewidth}{1.5pt}}\par\vspace{3pt}
\noindent{\fontsize{26}{28}\selectfont\bfseries\color{accent}11}\hspace{6pt}
{\LARGE\bfseries 介绍 · 告知信}\hspace{7pt}{\large\color{gray} Introduction}\hfill
{\footnotesize\color{gray} 真题 2020 · 2022 · 2024}\par
\vspace{1pt}
{\color{accent!35}\rule{\linewidth}{0.5pt}}\par\vspace{3pt}
{\footnotesize {\color{accent}\bfseries 题干信号}\quad 题干出现 tell sb. about / introduce / 介绍景点、计划、活动}\par
{\footnotesize {\color{accent}\bfseries 写作身份}\quad 组织者或知情者；对国际学生要放慢语速、少用难词}\par
{\footnotesize {\color{accent}\bfseries 必须写到}\quad ①介绍对象　②两三个特点　③一处亮点　④参观或参与提示}\par
{\footnotesize {\color{accent}\bfseries 一句话策略}\quad 第二段讲“是什么”，第三段讲“注意什么”，两段分工清楚}\par

\vspace{5pt}
{\color{accent}\bfseries 一、整篇模板（背这一篇）}\hfill{\footnotesize\color{gray} 正文 104 词 · 9 句}\par
\vspace{2pt}
\noindent\fcolorbox{rule}{soft}{\begin{minipage}{\dimexpr\linewidth-2\fboxsep-2\fboxrule\relax}\vspace{3pt}\small\linespread{1.03}\selectfont Dear {\color{hot}\bfseries 【称呼】},\par
\vspace{3pt}
{\color{accent}\bfseries ①}I am glad to tell you about {\color{hot}\bfseries 【介绍对象】}, which we are going to visit {\color{hot}\bfseries 【日期】}.\par
\vspace{2pt}
{\color{accent}\bfseries ②}It is {\color{hot}\bfseries 【特点一】}.{\color{accent}\bfseries ③}It is also {\color{hot}\bfseries 【特点二】}.{\color{accent}\bfseries ④}The part I like best is {\color{hot}\bfseries 【亮点】}.{\color{accent}\bfseries ⑤}I am sure you will find it interesting.\par
\vspace{2pt}
{\color{accent}\bfseries ⑥}Before we set off, I would like to give you two tips.{\color{accent}\bfseries ⑦}First, {\color{hot}\bfseries 【贴士一】}.{\color{accent}\bfseries ⑧}Second, {\color{hot}\bfseries 【贴士二】}.{\color{accent}\bfseries ⑨}If you have any questions, please let me know.\par
\vspace{2pt}
\vspace{2pt}Yours sincerely,\par
Li Ming\par\vspace{3pt}\end{minipage}}\par
\vspace{1pt}
{\footnotesize\color{gray} 段落分工｜① 点题｜②—⑤ 特点＋亮点｜⑥—⑨ 两条贴士}\par

\vspace{6pt}
{\color{accent}\bfseries 二、真题落地（2020 真题 · 介绍古镇并给参观提示）}\par
\vspace{2pt}
\noindent\fcolorbox{sandline}{sand}{\begin{minipage}{\dimexpr\linewidth-2\fboxsep-2\fboxrule\relax}\vspace{3pt}\small\linespread{1.03}\selectfont Dear friends,\par
\vspace{3pt}
{\color{accent}\bfseries ①}I am glad to tell you about Zhouzhuang, which we are going to visit next Saturday.\par
\vspace{2pt}
{\color{accent}\bfseries ②}It is famous for its waterways and stone bridges. {\color{accent}\bfseries ③}It is also home to more than nine hundred years of history. {\color{accent}\bfseries ④}The part I like best is the old opera stage, where local performers sing every afternoon. {\color{accent}\bfseries ⑤}I am sure you will find it interesting.\par
\vspace{2pt}
{\color{accent}\bfseries ⑥}Before we set off, I would like to give you two tips. {\color{accent}\bfseries ⑦}First, wear comfortable shoes, as we will walk a lot. {\color{accent}\bfseries ⑧}Second, bring some cash, for small shops there do not accept cards. {\color{accent}\bfseries ⑨}If you have any questions, please let me know.\par
\vspace{2pt}
\vspace{2pt}Yours sincerely,\par
Li Ming\par\vspace{3pt}\end{minipage}}\par

\vspace{6pt}
{\color{accent}\bfseries 三、换装表（只改这几处）}\par
\vspace{2pt}
{\footnotesize
\begin{tabularx}{\linewidth}{@{}>{\raggedright\arraybackslash}p{3.5cm}>{\raggedright\arraybackslash}p{5.6cm}X@{}}
\toprule
{\color{accent}\bfseries 槽位} & {\color{accent}\bfseries 本篇真题怎么填} & {\color{accent}\bfseries 可换词块} \\
\midrule
【介绍对象】 & Zhouzhuang & the Palace Museum / the campus library \\
【特点一】【特点二】 & waterways / nine hundred years & free entry / open until 9 p.m. \\
【贴士一】【贴士二】 & comfortable shoes / bring cash & book online / bring your student card \\
\bottomrule
\end{tabularx}}\par

\vspace{5pt}
\noindent\begin{minipage}[t]{0.555\linewidth}
{\color{accent}\bfseries 四、阅卷人找什么}\par
{\footnotesize\begin{itemize}[leftmargin=1.2em,itemsep=1pt,topsep=1pt,parsep=0pt]\setlength{\leftmargin}{1.1em}
\item 介绍对象有“身份”（古镇、博物馆、项目），不是 it is very good
\item 两条提示都是可执行的动作（穿什么、带什么、几点到）
\item 对国际学生避免生僻词，句子短而清楚
\end{itemize}}
\end{minipage}\hfill
\begin{minipage}[t]{0.415\linewidth}
{\color{hot}\bfseries 别踩这两个坑}\par
{\footnotesize\begin{itemize}[leftmargin=1.2em,itemsep=1pt,topsep=1pt,parsep=0pt]\setlength{\leftmargin}{1.1em}
\item 把介绍写成景点广告，全是形容词，没有一条实用信息
\item 提示写成 Again, it is beautiful，等于没写
\end{itemize}}
\end{minipage}\par

\vspace{6pt}
{\color{accent}\bfseries 五、想上 9 分，从这两句里挑一句加在结尾}\par
\vspace{2pt}
{\footnotesize
\begin{itemize}[leftmargin=1.2em,itemsep=1.5pt,topsep=1pt,parsep=0pt]
\item You will find the trip both relaxing and rewarding.\quad{\color{gray}（both…and 并列收束，比 very nice 具体）}
\item Do remember to charge your phone before we set off.\quad{\color{gray}（do 强调提醒，贴士更像贴士）}
\end{itemize}}\par

\vspace{6pt}
{\color{accent}\bfseries 六、题干变体怎么改（同一篇，换两个词）}\par
\vspace{2pt}
{\footnotesize
\begin{itemize}[leftmargin=1.2em,itemsep=1.5pt,topsep=1pt,parsep=0pt]
\item 2022 美食节：介绍对象换成 the festival，贴士改成空腹来、带现金。
\item 2025 短剧：介绍对象换成 the play，贴士改成提前入场、保持安静。
\end{itemize}}\par


\clearpage
\phantomsection\label{g12}
\noindent{\color{accent}\rule{\linewidth}{1.5pt}}\par\vspace{3pt}
\noindent{\fontsize{26}{28}\selectfont\bfseries\color{accent}12}\hspace{6pt}
{\LARGE\bfseries 推荐信}\hspace{7pt}{\large\color{gray} Recommendation}\hfill
{\footnotesize\color{gray} 真题 模拟高频（大纲内）}\par
\vspace{1pt}
{\color{accent!35}\rule{\linewidth}{0.5pt}}\par\vspace{3pt}
{\footnotesize {\color{accent}\bfseries 题干信号}\quad 题干出现 recommend / 推荐一本书、一个地方、一个人}\par
{\footnotesize {\color{accent}\bfseries 写作身份}\quad 了解被推荐对象的同学、朋友或同事}\par
{\footnotesize {\color{accent}\bfseries 必须写到}\quad ①推荐什么　②两三条理由　③一个具体事例　④获取方式或联系方式}\par
{\footnotesize {\color{accent}\bfseries 一句话策略}\quad 理由要“可验证”：内容、价格、使用体验；不要只写 highly recommend}\par

\vspace{5pt}
{\color{accent}\bfseries 一、整篇模板（背这一篇）}\hfill{\footnotesize\color{gray} 正文 99 词 · 9 句}\par
\vspace{2pt}
\noindent\fcolorbox{rule}{soft}{\begin{minipage}{\dimexpr\linewidth-2\fboxsep-2\fboxrule\relax}\vspace{3pt}\small\linespread{1.03}\selectfont Dear {\color{hot}\bfseries 【称呼】},\par
\vspace{3pt}
{\color{accent}\bfseries ①}I hear that {\color{hot}\bfseries 【推荐缘由】}.{\color{accent}\bfseries ②}I would like to recommend {\color{hot}\bfseries 【推荐对象】} without hesitation.\par
\vspace{2pt}
{\color{accent}\bfseries ③}{\color{hot}\bfseries 【对象】} is {\color{hot}\bfseries 【优点一】}.{\color{accent}\bfseries ④}Besides, {\color{hot}\bfseries 【优点二】}.{\color{accent}\bfseries ⑤}In my experience, {\color{hot}\bfseries 【具体事例】}.{\color{accent}\bfseries ⑥}I believe {\color{hot}\bfseries 【结论】}.\par
\vspace{2pt}
{\color{accent}\bfseries ⑦}It is available in most bookshops for less than fifty yuan.{\color{accent}\bfseries ⑧}If you need more information, please contact me.{\color{accent}\bfseries ⑨}I am sure you will not regret the choice.\par
\vspace{2pt}
\vspace{2pt}Yours sincerely,\par
Li Ming\par\vspace{3pt}\end{minipage}}\par
\vspace{1pt}
{\footnotesize\color{gray} 段落分工｜①—② 推荐什么｜③—⑥ 理由＋事例｜⑦—⑨ 获取方式}\par

\vspace{6pt}
{\color{accent}\bfseries 二、真题落地（模拟题 · 向朋友推荐一本中国史读物）}\par
\vspace{2pt}
\noindent\fcolorbox{sandline}{sand}{\begin{minipage}{\dimexpr\linewidth-2\fboxsep-2\fboxrule\relax}\vspace{3pt}\small\linespread{1.03}\selectfont Dear Jack,\par
\vspace{3pt}
{\color{accent}\bfseries ①}I hear that you are looking for a book on Chinese history. {\color{accent}\bfseries ②}I would like to recommend A Brief History of China without hesitation.\par
\vspace{2pt}
{\color{accent}\bfseries ③}The book is a clear guide to five thousand years of Chinese history. {\color{accent}\bfseries ④}Besides, each chapter ends with a short summary, which makes review easy. {\color{accent}\bfseries ⑤}In my experience, it helped me understand the dynasties far better than any other book. {\color{accent}\bfseries ⑥}I believe it will suit you too.\par
\vspace{2pt}
{\color{accent}\bfseries ⑦}It is available in most bookshops for less than fifty yuan. {\color{accent}\bfseries ⑧}If you need more information, please contact me. {\color{accent}\bfseries ⑨}I am sure you will not regret the choice.\par
\vspace{2pt}
\vspace{2pt}Yours sincerely,\par
Li Ming\par\vspace{3pt}\end{minipage}}\par

\vspace{6pt}
{\color{accent}\bfseries 三、换装表（只改这几处）}\par
\vspace{2pt}
{\footnotesize
\begin{tabularx}{\linewidth}{@{}>{\raggedright\arraybackslash}p{3.5cm}>{\raggedright\arraybackslash}p{5.6cm}X@{}}
\toprule
{\color{accent}\bfseries 槽位} & {\color{accent}\bfseries 本篇真题怎么填} & {\color{accent}\bfseries 可换词块} \\
\midrule
【推荐对象】 & A Brief History of China & the City Museum / Professor Wang \\
【优点一】【优点二】 & clear guide / chapter summaries & free entry / English guide available \\
【具体事例】 & it helped me understand the dynasties & I visited it twice last year \\
\bottomrule
\end{tabularx}}\par

\vspace{5pt}
\noindent\begin{minipage}[t]{0.555\linewidth}
{\color{accent}\bfseries 四、阅卷人找什么}\par
{\footnotesize\begin{itemize}[leftmargin=1.2em,itemsep=1pt,topsep=1pt,parsep=0pt]\setlength{\leftmargin}{1.1em}
\item 第一段一句话说清推荐什么，不铺垫
\item 理由两到三条，每条都能被验证（内容、价格、体验）
\item 有一个“我自己的使用经历”，可信度立刻上来
\end{itemize}}
\end{minipage}\hfill
\begin{minipage}[t]{0.415\linewidth}
{\color{hot}\bfseries 别踩这两个坑}\par
{\footnotesize\begin{itemize}[leftmargin=1.2em,itemsep=1pt,topsep=1pt,parsep=0pt]\setlength{\leftmargin}{1.1em}
\item 通篇 highly recommend / very good，没有一条具体理由
\item 推荐信写成广告，堆形容词，没有一句事实
\end{itemize}}
\end{minipage}\par

\vspace{6pt}
{\color{accent}\bfseries 五、想上 9 分，从这两句里挑一句加在结尾}\par
\vspace{2pt}
{\footnotesize
\begin{itemize}[leftmargin=1.2em,itemsep=1.5pt,topsep=1pt,parsep=0pt]
\item It is the kind of book you will finish in a weekend.\quad{\color{gray}（定语从句，评价具体可感）}
\item I would be glad to lend you my copy first.\quad{\color{gray}（主动提供帮助，结尾更暖）}
\end{itemize}}\par

\vspace{6pt}
{\color{accent}\bfseries 六、题干变体怎么改（同一篇，换两个词）}\par
\vspace{2pt}
{\footnotesize
\begin{itemize}[leftmargin=1.2em,itemsep=1.5pt,topsep=1pt,parsep=0pt]
\item 推荐人：优点换成 works hard / has a good command of English。
\item 推荐地点：优点换成 free entry / easy to reach，事例写自己去过。
\end{itemize}}\par


\clearpage
\vspace*{-4pt}
{\color{accent}\bfseries\LARGE 交卷前 60 秒：自评与自查}\par\vspace{3pt}
{\color{accent}\rule{\linewidth}{1.1pt}}\par
\vspace{6pt}

\noindent\begin{minipage}[t]{0.49\linewidth}\raggedright
{\color{accent}\bfseries 一、十二条打勾清单}\par
{\footnotesize\begin{itemize}[leftmargin=1.2em,itemsep=1.5pt,topsep=1pt,parsep=0pt]
\item[$\square$] 两个任务动作都写了吗（题干的 1) 2)）
\item[$\square$] 时间、地点、数字写了吗
\item[$\square$] 称呼有逗号，结束语有逗号，署名顶格吗
\item[$\square$] 通知有标题和右下角落款吗
\item[$\square$] 时态统一吗（邀请/安排用将来，感谢/投诉用过去）
\item[$\square$] 主谓一致吗（there are / it is / people are）
\item[$\square$] 拼写检查了吗（receive, convenient, sincerely）
\item[$\square$] 正文 95—105 词吗（不含称呼与落款）
\item[$\square$] 每段都有连接词吗（As for / Besides / Therefore）
\item[$\square$] 有没有写“I am a student”这类无关背景
\item[$\square$] 有没有把题目姓名写成自己的名字
\item[$\square$] 有没有空行分三段（通知可分两段）
\end{itemize}}
\end{minipage}\hfill
\begin{minipage}[t]{0.49\linewidth}\raggedright
{\color{accent}\bfseries 二、按档位给自己定位}\par
{\footnotesize
漏一个任务动作 → 最高 6 分档；漏两个 → 4 分档以下。\par\vspace{3pt}
格式错一处（称呼/落款/通知标题）→ 6—7 分档；错两处以上 → 5 分档。\par\vspace{3pt}
语言偶有小错但意思清楚 → 6—7 分；错误影响理解 → 4—5 分。\par\vspace{3pt}
要点全、格式对、有两处自然衔接、无明显错误 → 8—9 分。\par\vspace{3pt}
以上都对，且句式有变化、用词准确 → 可以到 10 分，不人为封顶。\par}
\end{minipage}\par

\vspace{12pt}
{\color{accent}\bfseries 三、十二组高频错误（左边一律改成右边）}\par
{\color{rule}\rule{\linewidth}{0.5pt}}\par\vspace{3pt}
{\footnotesize
\begin{tabularx}{\linewidth}{@{}>{\raggedright\arraybackslash}X>{\raggedright\arraybackslash}X@{}}
\toprule
{\color{hot}\bfseries 错误写法} & {\color{accent}\bfseries 正确写法} \\
\midrule
congratulations for & congratulations on \\
look forward to meet you & look forward to meeting you \\
I would appreciate if & I would appreciate it if \\
grateful to your help & grateful for your help \\
suggest you to go & suggest that you go \\
There will have a meeting & There will be a meeting \\
I look forward to hear from you & I look forward to hearing from you \\
a lot of informations & a lot of information \\
take part in the meeting & attend the meeting \\
I am agree with you & I agree with you \\
pay attention on & pay attention to \\
It is high time we start & It is high time we started \\
\bottomrule
\end{tabularx}}\par

\vspace{12pt}
{\color{accent}\bfseries 四、通用零件（不够用时从这里取，别再另背一套）}\par
{\color{rule}\rule{\linewidth}{0.5pt}}\par\vspace{4pt}
{\footnotesize
\begin{tabularx}{\linewidth}{@{}>{\raggedright\arraybackslash}p{2.1cm}X@{}}
\toprule
{\color{accent}\bfseries 功能} & {\color{accent}\bfseries 现成表达} \\
\midrule
开头 & I am writing to invite / thank / complain about…；I was delighted to learn that… \\
收尾 & I look forward to your reply. / Thank you for your time and consideration. \\
理由 & For one thing… For another…；The main reason is that… \\
请求 & I would be grateful if you could…；Would it be possible for you to…? \\
时间地点 & next Saturday at 3 p.m. in Room 203；from July 10 to July 20；by Friday \\
人数与费用 & twenty volunteers；free of charge；a deposit of 200 yuan \\
\bottomrule
\end{tabularx}}\par

\vspace{10pt}
\noindent\colorbox{sand}{\parbox{\dimexpr\linewidth-2\fboxsep\relax}{%
\vspace{4pt}{\footnotesize
{\bfseries 最后一句：}小作文不是比谁句子长，是比谁把两个动作写全、把格式写对、把话说清楚。\\
背完整篇之后，一定要掐 15 分钟在答题卡上默写一遍；模板只有落到纸上，才是你的。}
\vspace{4pt}}}

\end{document}